%% 3-11 指数関数：いまある量に比例してふえる・へる（第12回・組 V3-6）  % noruby
%% 流れ：倍々の表（やってみよう）→ ①ふえる量÷いまある量が一定 → ②y=a^x のグラフ → ③指数のきまり  % noruby
%%       → ④半減期 → ⑤分数の指数（表と (a^m)^n の2通り）→ ⑥整数でない指数を電卓で。e と ln は使わない（5-14）。  % noruby
\makeatletter
% 細菌1個（緑のカプセル）  \bkthrlcell{(x,y)}{向き（度）}  % noruby
\newcommand\bkthrlcell[2]{\begin{scope}[shift={#1}, rotate=#2]
  \filldraw[fill=bkGreen!28, draw=bkGreen!70!black, line width=0.6pt, rounded corners=1.1mm] (-2.3,-1.15) rectangle (2.3,1.15);
  \fill[bkGreen!70!black] (-0.8,0.15) circle (0.32) (0.9,-0.25) circle (0.28);
\end{scope}}
% おにぎり  \bkthrlonigiri{(x,y)}{大きさ（mm）}  % noruby
\newcommand\bkthrlonigiri[2]{\begin{scope}[shift={#1}, scale=#2/10]
  \filldraw[fill=white, draw=bkGray!80!black, line width=0.7pt, rounded corners=2.6mm] (-6,-4.5) -- (6,-4.5) -- (0,6) -- cycle;
  \fill[bkInk!85, rounded corners=0.5mm] (-2.6,-4.7) rectangle (2.6,-0.6);
  \foreach \bkthrl@p in {(-2.6,1.2),(1.9,2.2),(-0.4,3.6),(3.4,-1.4),(-3.9,-2)} {\fill[bkGray!40] \bkthrl@p circle (0.25);}
\end{scope}}
% かけ算の札（2-12 のテープと同じ形）  \bkthrltape{(x,y)}{はば}{色}{文字}  % noruby
\newcommand\bkthrltape[4]{\begin{scope}[shift={#1}]
  \filldraw[fill=#3!18, draw=#3!80!black, line width=0.7pt, rounded corners=1mm] (0,0) rectangle (#2,6.5);
  \node[font=\small] at ({#2/2},3.25) {#4};
\end{scope}}
% 説明の段の見出し（丸数字＋太字の1行）  \bkthrlstep{番号}{見出し}  % noruby
\newcommand\bkthrlstep[2]{\par\addvspace{6pt}\Needspace{5\baselineskip}\noindent
  {\bkcircled[bkOrange]{#1}\enspace\sffamily\bfseries\color{bkOrange!75!black}#2}\par\nobreak\vspace{2pt}}
% 式を、まん中に1行で置く  \bkthrlline{…}  % noruby
\newcommand\bkthrlline[1]{\par\addvspace{3pt}{\centering #1\par}\addvspace{3pt}}
% 表計算で目立たせるセルの色  % noruby
\colorlet{bkthrlsel}{bkOrange!22}
\makeatother

\bkchapter{3}{11}{\R{指数関数}{し|すう|かん|すう}：いまある\R{量}{りょう}に\\\R{比例}{ひ|れい}してふえる・へる}{\bklevelrange{2}{3}}{\R{指数関数}{し|すう|かん|すう}：いまある\R{量}{りょう}に\R{比例}{ひ|れい}してふえる・へる}
\bkmichishirube{2-7,2-12,2-18,3-3,3-5}{3-12,5-14,5-18,6-4}

\begin{goal}
\begin{itemize}[label={\color{\bkTT}●}, leftmargin=1.6\zw, itemsep=0pt, topsep=1pt]
\item \bklevel{2}\enspace 「ふえる\R{量}{りょう}が、いまある\R{量}{りょう}に\R{比例}{ひ|れい}する」\R{変化}{へん|か}を\R{表}{ひょう}にして、\R{指数関数}{し|すう|かん|すう}の\R{式}{しき}とグラフに\R{表}{あらわ}せる。
\item \bklevel{2}\enspace \R{指数}{し|すう}のきまり $a^{x+h}=a^x\times a^h$、$(a^m)^n=a^{mn}$、$a^0=1$ を\R{使}{つか}って\R{計算}{けい|さん}できる。
\item \bklevel{3}\enspace $x^{\frac{1}{2}}=\sqrt{x}$ のわけがわかり、\R{半減期}{はん|げん|き}や\R{整数}{せい|すう}でない\R{指数}{し|すう}を\R{電卓}{でん|たく}で\R{計算}{けい|さん}できる。
\end{itemize}
\end{goal}

\begin{hatena}
\noindent\begin{minipage}[c]{0.53\linewidth}
\R{朝}{あさ}8\R{時}{じ}に\R{作}{つく}ったおにぎりに、\R{細菌}{さい|きん}が1\R{個}{こ}ついていました。この\R{細菌}{さい|きん}は、あたたかい\R{所}{ところ}では、20\R{分}{ぷん}ごとに1\R{個}{こ}が2\R{個}{こ}に\R{分}{わ}かれます。

お\R{昼}{ひる}の12\R{時}{じ}には、\R{細菌}{さい|きん}は\R{何個}{なん|こ}になっているかな？　\R{夕方}{ゆう|がた}の6\R{時}{じ}には？　1\R{個}{こ}が2\R{個}{こ}になるだけなら、たいしたことはなさそうだけれど……。
\end{minipage}\hfill
\begin{minipage}[c]{0.44\linewidth}\centering
\begin{tikzpicture}[x=1.0mm, y=1.0mm]
  \bkthrlonigiri{(4,11)}{13}
  \bkthrlcell{(4,-6)}{0}
  \bkthrlcell{(21,-3)}{20}\bkthrlcell{(21,-9)}{-15}
  \bkthrlcell{(38,0)}{10}\bkthrlcell{(38,-4.5)}{-20}\bkthrlcell{(38,-9)}{25}\bkthrlcell{(38,-13.5)}{0}
  \node[font=\sffamily\bfseries\Large, text=bkOrange!80!black] at (47,-6.5) {？};
  \foreach \a/\b in {8/16, 25/33} {\draw[bk arrow, bkGray, line width=0.7pt] (\a,-6) -- (\b,-6);}
  \node[font=\sffamily\scriptsize, anchor=north] at (4,-9) {1\R{個}{こ}};
  \node[font=\sffamily\scriptsize, anchor=north] at (21,-12.5) {2\R{個}{こ}};
  \node[font=\sffamily\scriptsize, anchor=north] at (38,-16.5) {4\R{個}{こ}};
  \node[font=\sffamily\scriptsize, anchor=south] at (12,-4.2) {20\R{分}{ぷん}};
  \node[font=\sffamily\scriptsize, anchor=south] at (29,-4.2) {20\R{分}{ぷん}};
\end{tikzpicture}
\end{minipage}
\end{hatena}

ふえ\R{方}{かた}には、「\R{同}{おな}じ\R{数}{かず}ずつ」（3-5の1\R{次関数}{じ|かん|すう}）と「\R{同}{おな}じ\R{倍}{ばい}ずつ」があります。この\R{章}{しょう}では、「\R{同}{おな}じ\R{倍}{ばい}ずつ」ふえる・へる\R{量}{りょう}を\R{調}{しら}べます。\R{細菌}{さい|きん}、お\R{湯}{ゆ}の\R{冷}{さ}め\R{方}{かた}、\R{放射線}{ほう|しゃ|せん}が、このなかまです。

%% ---------------- やってみよう・思い出そう ----------------
\begin{yattemiyou}[「100\R{個}{こ}ずつ」と「2\R{倍}{ばい}ずつ」、どっちが\R{多}{おお}い？]
\R{表計算}{ひょう|けい|さん}で、2つのふえ\R{方}{かた}をくらべます（2-18）。どちらも、はじめは1\R{個}{こ}です。B\R{列}{れつ}は「20\R{分}{ぷん}ごとに2\R{倍}{ばい}」、C\R{列}{れつ}は「20\R{分}{ぷん}ごとに100\R{個}{こ}ずつふえる」です。
\begin{enumerate}[label={\bkcircled[bkGreen]{\arabic*}}, leftmargin=2.4\zw, itemsep=1pt, topsep=3pt]
\item 2\R{行目}{ぎょう|め}に、はじめのようすを\R{入}{い}れよう：A2に0、B2に1、C2に1。
\item 3\R{行目}{ぎょう|め}に、1つ\R{上}{うえ}の\R{行}{ぎょう}を\R{使}{つか}う\R{式}{しき}を\R{入}{い}れよう：A3に \texttt{=A2+1}、B3に \texttt{=B2*2}、C3に \texttt{=C2+100}。
\item 3\R{行目}{ぎょう|め}を\R{選}{えら}んで、14\R{行目}{ぎょう|め}まで\R{下}{した}へコピーしよう。
\item B\R{列}{れつ}とC\R{列}{れつ}をくらべよう。B\R{列}{れつ}がC\R{列}{れつ}をぬくのは\R{何回目}{なん|かい|め}かな？
\end{enumerate}
\par\vspace{2pt}
\noindent\begin{minipage}[t]{0.47\linewidth}\vspace{0pt}\centering
\bksheet{3}{14}{%
1 & \R{回}{かい} & 2\R{倍}{ばい}ずつ & 100ずつ \\
2 & 0 & 1 & 1 \\
3 & \texttt{=A2+1} & \SetCell{bg=bkthrlsel}\texttt{=B2*2} & \texttt{=C2+100} \\
4 & \texttt{=A3+1} & \texttt{=B3*2} & \texttt{=C3+100} \\}
\end{minipage}\hfill
\begin{minipage}[t]{0.5\linewidth}\vspace{0pt}
A\R{列}{れつ}の「\R{回}{かい}」は、20\R{分}{ぷん}が\R{何回}{なん|かい}たったかです。1\R{回}{かい}で20\R{分}{ぷん}、3\R{回}{かい}で60\R{分}{ぷん}（1\R{時間}{じ|かん}）です。\texttt{*} はかけ\R{算}{ざん}の\R{記号}{き|ごう}です（2-18）。
\end{minipage}
\end{yattemiyou}

\noindent \R{結果}{けっ|か}は\R{次}{つぎ}のようになります。\R{上}{うえ}の\R{行}{ぎょう}の\R{答}{こた}えを\R{使}{つか}って、\R{下}{した}の\R{行}{ぎょう}を\R{計算}{けい|さん}しています（2-18のリレー）。

\begin{bkcenter}
\begin{bktable}{colspec={Q[l,m,wd=26mm] *{7}{Q[c,m,wd=11mm]}}, rowsep=1.2pt, cells={font=\small}, column{1}={font=\sffamily\small}, row{1,4}={bg=bkTheme!10}, cell{5}{5-7}={bg=bkthrlsel}, hline{4}={1pt,bkLine}}
\R{回}{かい}（20\R{分}{ぷん}ごと） & 0 & 1 & 2 & 3 & 4 & 5 & 6 \\
B：2\R{倍}{ばい}ずつ & 1 & 2 & 4 & 8 & 16 & 32 & 64 \\
C：100ずつ & 1 & 101 & 201 & 301 & 401 & 501 & 601 \\
\R{回}{かい}（20\R{分}{ぷん}ごと） & 7 & 8 & 9 & 10 & 11 & 12 & \\
B：2\R{倍}{ばい}ずつ & 128 & 256 & 512 & 1024 & 2048 & 4096 & \\
C：100ずつ & 701 & 801 & 901 & 1001 & 1101 & 1201 & \\
\end{bktable}
\bktabcaption{10\R{回目}{かい|め}（200\R{分後}{ぷん|ご}）で、B\R{列}{れつ}がC\R{列}{れつ}をぬく。}
\end{bkcenter}

\noindent はじめのうちは、「100\R{個}{こ}ずつ」のほうが、ずっと\R{多}{おお}くふえます。6\R{回目}{かい|め}でも、Bは64\R{個}{こ}、Cは601\R{個}{こ}です。ところが、10\R{回目}{かい|め}でBは1024\R{個}{こ}になり、Cの1001\R{個}{こ}をぬきます。そこから\R{先}{さき}は、\R{差}{さ}がどんどん\R{開}{ひら}きます。12\R{回目}{かい|め}には、Bが4096\R{個}{こ}、Cは1201\R{個}{こ}です。

\begin{omoidasou}
\bkrecall{2-7}{\R{同}{おな}じ\R{数}{かず}を\R{何回}{なん|かい}もかけるのが\R{累乗}{るい|じょう}。$2\times2\times2=2^3=8$。$2^{10}=1024$ は、ほぼ1000。}
\bkrecall{2-12}{$x^m\times x^n=x^{m+n}$（かけ\R{算}{ざん}は\R{指数}{し|すう}をたす）。$x^0=1$、\ $x^{-1}=\dfrac{1}{x}$、\ $x^{-2}=\dfrac{1}{x^2}$}
\bkrecall{2-5}{10\,\%ふえる＝もとの1\R{倍}{ばい}＋0.1\R{倍}{ばい}＝1.1\R{倍}{ばい}。20\,\%へる＝0.8\R{倍}{ばい}。}
\bkrecall{2-8}{$\sqrt{2}$ は、2\R{乗}{じょう}すると2になる\R{正}{せい}の\R{数}{かず}。$\sqrt{2}=1.414\cdots$、\ $\sqrt{9}=3$}
\bkrecall{3-5}{1\R{次関数}{じ|かん|すう} $y=ax+b$ は、$x$ が1ふえるごとに、$y$ がいつも\R{同}{おな}じ $a$ だけふえる。}
\end{omoidasou}

\bkfukidashi{sora}{はじめは「100\R{個}{こ}ずつ」の\R{楽勝}{らく|しょう}だと\R{思}{おも}ったのに、\R{途中}{と|ちゅう}で\R{一気}{いっ|き}にぬかれた！　\R{倍々}{ばい|ばい}って、どうしてこんなに\R{強}{つよ}いの？}

\clearpage
%% ---------------- 数学の道具 ----------------
\bksection{doguu}{\R{数学}{すう|がく}の\R{道具}{どう|ぐ}}

\bksubsection{1}{ふえる\R{量}{りょう}が、いまある\R{量}{りょう}に\R{比例}{ひ|れい}する}

\leavevmode\label{bk:3-11:hirei}やってみようの\R{表}{ひょう}で、1\R{回}{かい}ごとに「いくつふえたか」を\R{調}{しら}べます。ふえた\R{数}{かず}は、あと－まえ（2-3）で\R{求}{もと}めます。

\begin{bkcenter}
\begin{bktable}{colspec={Q[l,m,wd=41mm] *{6}{Q[c,m,wd=11mm]}}, rowsep=1.2pt, cells={font=\small}, column{1}={font=\sffamily\small}}
\R{回}{かい} & 0 & 1 & 2 & 3 & 4 & 5 \\
C：100ずつ & 1 & 101 & 201 & 301 & 401 & 501 \\
\SetCell{bg=bkBlue!10} C　ふえた\R{数}{かず} & & 100 & 100 & 100 & 100 & 100 \\
B：2\R{倍}{ばい}ずつ & 1 & 2 & 4 & 8 & 16 & 32 \\
\SetCell{bg=bkthrlsel} B　ふえた\R{数}{かず} & & 1 & 2 & 4 & 8 & 16 \\
\end{bktable}
\bktabcaption{Cは\R{毎回}{まい|かい}100ずつ。Bのふえた\R{数}{かず}は、\R{毎回}{まい|かい}ちがう。}
\end{bkcenter}

Cは、\R{毎回}{まい|かい}\R{同}{おな}じ100ずつふえています。3-5で\R{学}{まな}んだ1\R{次関数}{じ|かん|すう}と\R{同}{おな}じで、\R{変化}{へん|か}の\R{割合}{わり|あい}がいつも\R{同}{おな}じです。

Bのふえた\R{数}{かず}は、1、2、4、8、16と、\R{毎回}{まい|かい}ちがいます。でも、よく\R{見}{み}ると、ふえた\R{数}{かず}は、ふえる\R{前}{まえ}にあった\R{数}{かず}と\R{同}{おな}じです。たとえば、3\R{回目}{かい|め}の8\R{個}{こ}から4\R{回目}{かい|め}の16\R{個}{こ}へは、8\R{個}{こ}ふえています。8\R{個}{こ}の\R{細菌}{さい|きん}が、それぞれ2つに\R{分}{わ}かれたからです。

\begin{bkfigure}
\begin{tikzpicture}[x=1mm, y=1mm]
  % 4回目の16個（8個がそれぞれ分かれる）  % noruby
  \foreach \i in {0,...,7} {\bkthrlcell{(8+\i*7,10)}{0}}
  \node[font=\sffamily\footnotesize, anchor=east] at (2,10) {3\R{回目}{かい|め}};
  \node[font=\sffamily\footnotesize, anchor=west] at (61,10) {8\R{個}{こ}};
  \foreach \i in {0,...,7} {\draw[bkGray, line width=0.5pt] (8+\i*7,8.2) -- (6.6+\i*7,3.6) (8+\i*7,8.2) -- (9.4+\i*7,3.6);
    \bkthrlcell{(6.4+\i*7,1)}{80}
    \begin{scope}[shift={(9.6+\i*7,1)}, rotate=80]
      \filldraw[fill=bkOrange!35, draw=bkOrange!80!black, line width=0.6pt, rounded corners=1.1mm] (-2.3,-1.15) rectangle (2.3,1.15);
    \end{scope}}
  \node[font=\sffamily\footnotesize, anchor=east] at (2,1) {4\R{回目}{かい|め}};
  \node[font=\sffamily\footnotesize, anchor=west, align=left] at (61,1) {16\R{個}{こ}};
  \draw[decorate, decoration={brace, amplitude=2.5pt, mirror}, line width=0.6pt, bkOrange!80!black] (5,-3.5) -- (61,-3.5);
  \node[font=\sffamily\footnotesize, text=bkOrange!80!black, anchor=north] at (33,-6) {ふえた\R{分}{ぶん}（オレンジ）は8\R{個}{こ}＝ふえる\R{前}{まえ}にあった\R{数}{かず}};
  % 右：棒の図  % noruby
  \begin{scope}[shift={(84,-6)}]
    \foreach \k/\h in {0/1,1/2,2/4,3/8,4/16} {
      \pgfmathsetmacro\hh{\h*1.45}\pgfmathsetmacro\hp{\h*1.45/2}
      \fill[bkGreen!40] (\k*8,0) rectangle (\k*8+5.5,\hp);
      \ifnum\k>0 \fill[bkOrange!45] (\k*8,\hp) rectangle (\k*8+5.5,\hh);\else \fill[bkGreen!40] (0,0) rectangle (5.5,\hh);\fi
      \draw[bkInk!70, line width=0.5pt] (\k*8,0) rectangle (\k*8+5.5,\hh);
      \node[font=\sffamily\scriptsize, anchor=north] at (\k*8+2.75,-0.8) {\k};}
    \draw[bkInk, line width=0.7pt] (-1.5,0) -- (39,0);
    \node[font=\sffamily\scriptsize, anchor=north west] at (37,-0.8) {\R{回}{かい}};
  \end{scope}
\end{tikzpicture}
\bkfigcaption{\R{左}{ひだり}：8\R{個}{こ}がそれぞれ2つに\R{分}{わ}かれると、8\R{個}{こ}ふえる。\R{右}{みぎ}：\R{棒}{ぼう}のオレンジの\R{部分}{ぶ|ぶん}（ふえた\R{分}{ぶん}）は、\R{前}{まえ}の\R{棒}{ぼう}と\R{同}{おな}じ\R{高}{たか}さ。}
\end{bkfigure}

そこで、「ふえた\R{数}{かず}÷ふえる\R{前}{まえ}の\R{数}{かず}」を\R{計算}{けい|さん}してみます。2-15で、$y\div x$ がいつも\R{同}{おな}じなら、$y$ は $x$ に\R{比例}{ひ|れい}すると\R{学}{まな}びました。

\begin{bkcenter}
\begin{bktable}{colspec={Q[l,m,wd=48mm] *{5}{Q[c,m,wd=11mm]}}, rowsep=1.2pt, cells={font=\small}, column{1}={font=\sffamily\small}}
ふえる\R{前}{まえ}の\R{数}{かず}（いまある\R{量}{りょう}） & 1 & 2 & 4 & 8 & 16 \\
ふえた\R{数}{かず}（ふえる\R{量}{りょう}） & 1 & 2 & 4 & 8 & 16 \\
\SetCell{bg=bkthrlsel} ふえる\R{量}{りょう}÷いまある\R{量}{りょう} & 1 & 1 & 1 & 1 & 1 \\
\end{bktable}
\bktabcaption{わり\R{算}{ざん}の\R{答}{こた}えがいつも1。ふえる\R{量}{りょう}は、いまある\R{量}{りょう}に\R{比例}{ひ|れい}している。}
\end{bkcenter}

わり\R{算}{ざん}の\R{答}{こた}えは、いつも1です。つまり、\textbf{ふえる\R{量}{りょう}は、いまある\R{量}{りょう}に\R{比例}{ひ|れい}しています}（\R{比例定数}{ひ|れい|てい|すう}は1）。たくさんあるほど、たくさんふえます。だから、\R{倍々}{ばい|ばい}のふえ\R{方}{かた}は、\R{数}{かず}が\R{大}{おお}きくなるほど、ふえ\R{方}{かた}が\R{速}{はや}くなるのです。やってみようでBがCをぬいたのは、このためです。

\bkthrlstep{1}{もう1つの\R{例}{れい}：\R{利息}{り|そく}がつく\R{貯金}{ちょ|きん}}
1\R{年}{ねん}ごとに、いまある\R{金額}{きん|がく}の10\,\%の\R{利息}{り|そく}（ふえる\R{分}{ぶん}）がつく\R{貯金}{ちょ|きん}に、10000\R{円}{えん}を\R{預}{あず}けたとします。\R{利息}{り|そく}は、いまある\R{金額}{きん|がく}の0.1\R{倍}{ばい}です（10\,\%＝0.1。2-5）。

\begin{bkcenter}
\begin{bktable}{colspec={Q[l,m,wd=48mm] *{4}{Q[c,m,wd=15mm]}}, rowsep=1.2pt, cells={font=\small}, column{1}={font=\sffamily\small}}
\R{年}{ねん} & 0 & 1 & 2 & 3 \\
いまある\R{金額}{きん|がく}（\R{円}{えん}） & 10000 & 11000 & 12100 & 13310 \\
その\R{年}{とし}につく\R{利息}{り|そく}（\R{円}{えん}） & 1000 & 1100 & 1210 & \\
\SetCell{bg=bkthrlsel} \R{利息}{り|そく}÷いまある\R{金額}{きん|がく} & 0.1 & 0.1 & 0.1 & \\
\end{bktable}
\bktabcaption{\R{利息}{り|そく}は、いまある\R{金額}{きん|がく}に\R{比例}{ひ|れい}する（\R{比例定数}{ひ|れい|てい|すう}0.1）。}
\end{bkcenter}

1\R{年目}{ねん|め}の\R{利息}{り|そく}は、$10000\times0.1=1000$\R{円}{えん}です。2\R{年目}{ねん|め}は、11000\R{円}{えん}に\R{利息}{り|そく}がつくので、$11000\times0.1=1100$\R{円}{えん}です。\R{利息}{り|そく}にも\R{利息}{り|そく}がつくので、\R{利息}{り|そく}は\R{毎年}{まい|とし}すこしずつ\R{大}{おお}きくなります。

いまある\R{金額}{きん|がく}を1\R{倍}{ばい}として、それに0.1\R{倍}{ばい}の\R{利息}{り|そく}をたすので、1\R{年}{ねん}たつと $1+0.1=1.1$ \R{倍}{ばい}になります（\R{図}{ず}\ref{bk:3-11:fig-bai}）。
\begin{bkfigure}
\begin{tikzpicture}[x=1mm, y=1mm]
  \fill[bkBlue!25] (0,0) rectangle (12,20); \draw[bkBlue!80!black, line width=0.6pt] (0,0) rectangle (12,20);
  \fill[bkBlue!25] (44,0) rectangle (56,20); \fill[bkOrange!45] (44,20) rectangle (56,22);
  \draw[bkBlue!80!black, line width=0.6pt] (44,0) rectangle (56,20); \draw[bkOrange!80!black, line width=0.6pt] (44,20) rectangle (56,22);
  \draw[bk arrow, bkGray, line width=0.7pt] (15,10) -- (41,10);
  \node[font=\sffamily\footnotesize, anchor=south] at (28,11) {1\R{年}{ねん}たつと};
  \node[font=\sffamily\footnotesize, anchor=north] at (6,-1.5) {いま（10000\R{円}{えん}）};
  \node[font=\sffamily\footnotesize, anchor=north] at (50,-1.5) {1\R{年後}{ねん|ご}（11000\R{円}{えん}）};
  \node[font=\sffamily\footnotesize] at (6,10) {1\R{倍}{ばい}};
  \node[font=\sffamily\footnotesize] at (50,10) {1\R{倍}{ばい}};
  \draw[bkOrange!80!black, line width=0.5pt] (56.5,21) -- (59.5,21);
  \node[font=\sffamily\footnotesize, text=bkOrange!80!black, anchor=west] at (60,21) {＋0.1\R{倍}{ばい}（\R{利息}{り|そく}1000\R{円}{えん}）};
  \node[font=\sffamily\footnotesize, anchor=south] at (50,24) {あわせて1.1\R{倍}{ばい}};
\end{tikzpicture}
\bkfigcaption{いまある\R{金額}{きん|がく}（1\R{倍}{ばい}）に、その0.1\R{倍}{ばい}の\R{利息}{り|そく}がつくと、1.1\R{倍}{ばい}になる。}\label{bk:3-11:fig-bai}
\end{bkfigure}

だから、\R{表}{ひょう}は\R{次}{つぎ}のように、1.1を\R{何回}{なん|かい}もかけて\R{作}{つく}れます。
\[ 10000\times1.1=11000,\qquad 11000\times1.1=12100,\qquad 12100\times1.1=13310 \]
3\R{年後}{ねん|ご}は、10000に1.1を3\R{回}{かい}かけたので、$10000\times1.1^3=13310$\R{円}{えん}です（2-7の\R{累乗}{るい|じょう}）。

\bkthrlstep{2}{きまりにする}
2つの\R{例}{れい}を\R{並}{なら}べると、\R{同}{おな}じしくみが\R{見}{み}えます。
\begin{bkcenter}
\begin{bktable}{colspec={Q[l,m,wd=34mm] Q[c,m,wd=29mm] Q[c,m,wd=31mm] Q[c,m,wd=26mm]}, rowsep=1.5pt, cells={font=\small}, row{1}={font=\sffamily\bfseries\small, bg=bkTheme!10}}
 & ふえる\R{量}{りょう}は、いまある\R{量}{りょう}の & 1\R{回}{かい}ごとに & $x$ \R{回後}{かい|ご} \\
\R{細菌}{さい|きん}（はじめ1） & 1\R{倍}{ばい} & $1+1=2$ \R{倍}{ばい} & $1\times2^x$ \\
\R{貯金}{ちょ|きん}（はじめ10000） & 0.1\R{倍}{ばい} & $1+0.1=1.1$ \R{倍}{ばい} & $10000\times1.1^x$ \\
\end{bktable}
\bktabcaption{ふえる\R{量}{りょう}がいまある\R{量}{りょう}の $k$ \R{倍}{ばい}なら、1\R{回}{かい}ごとに $(1+k)$ \R{倍}{ばい}になる。}
\end{bkcenter}
1\R{回}{かい}ごとにかける\R{数}{かず}（\R{倍率}{ばい|りつ}）を $a$ とすると、
\begin{enumerate}[label={\bkcircled[bkTheme]{\arabic*}}, leftmargin=2.4\zw, itemsep=1pt, topsep=2pt]
\item 1\R{回}{かい}たつと、はじめの\R{量}{りょう}の $a$ \R{倍}{ばい}。
\item 2\R{回}{かい}たつと、それのまた $a$ \R{倍}{ばい}なので、$a\times a=a^2$ \R{倍}{ばい}。
\item $x$ \R{回}{かい}たつと、$a$ を $x$ \R{回}{かい}かけるので、$a^x$ \R{倍}{ばい}。
\end{enumerate}
\bkthrlline{\fbox{（$x$ \R{回後}{かい|ご}の\R{量}{りょう}）＝（はじめの\R{量}{りょう}）×$a^x$}}

\noindent へるときも\R{同}{おな}じです。\R{毎回}{まい|かい}、いまある\R{量}{りょう}の0.2\R{倍}{ばい}がへるなら、1\R{回}{かい}ごとに $1-0.2=0.8$ \R{倍}{ばい}になります。へる\R{量}{りょう}が、いまある\R{量}{りょう}に\R{比例}{ひ|れい}しているのです。

\bkthrlstep{3}{はてな？の\R{答}{こた}え}
\bkhatenaref の\R{細菌}{さい|きん}は、20\R{分}{ぷん}ごとに2\R{倍}{ばい}です。8\R{時}{じ}から12\R{時}{じ}までは\qty{4}{h}です。\R{分}{ふん}に\R{直}{なお}すと（2-10）、
\[ \qty{4}{h}\times\qty{60}{min/h}=\qty{240}{min} \]
20\R{分}{ぷん}が $240\div20=12$ \R{回}{かい}あるので、$1\times2^{12}=4096$\R{個}{こ}です。

\R{夕方}{ゆう|がた}6\R{時}{じ}（18\R{時}{じ}）までは\qty{10}{h}、つまり\qty{600}{min}です。$600\div20=30$ \R{回}{かい}なので、
\[ 1\times2^{30}=1073741824 \]
\R{約}{やく}10\R{億}{おく}7000\R{万}{まん}\R{個}{こ}です。\R{日本}{に|ほん}の\R{人口}{じん|こう}（\R{約}{やく}1\R{億}{おく}2000\R{万}{まん}\R{人}{にん}）の\R{約}{やく}9\R{倍}{ばい}です。\R{実際}{じっ|さい}には、\R{栄養}{えい|よう}が\R{足}{た}りなくなったり、\R{温度}{おん|ど}が\R{低}{ひく}かったりして、いつまでも\R{倍々}{ばい|ばい}ではふえません。それでも、\R{食}{た}べ\R{物}{もの}をあたたかい\R{所}{ところ}に\R{長}{なが}く\R{置}{お}いてはいけないわけが、\R{数}{かず}でよくわかります。

\bksubsection{2}{\R{指数関数}{し|すう|かん|すう} $y=a^x$ のグラフ}

\leavevmode\label{bk:3-11:graph}\R{細菌}{さい|きん}の\R{数}{かず}を $y$ \R{個}{こ}、20\R{分}{ぷん}が\R{何回}{なん|かい}たったかを $x$ とすると、はじめが1\R{個}{こ}なので $y=2^x$ です。$x$ を1つ\R{決}{き}めると $y$ が1つ\R{決}{き}まるので、$y$ は $x$ の\R{関数}{かん|すう}です（3-3）。

ここからは、$y$ を「$x=0$ のときの\R{量}{りょう}の\R{何倍}{なん|ばい}か」と\R{読}{よ}みます。すると、$x$ がマイナスのときも\R{考}{かんが}えられます。$x=-1$ は「1\R{回前}{かい|まえ}」、つまり20\R{分前}{ぷん|まえ}です。20\R{分}{ぷん}で2\R{倍}{ばい}になって $x=0$ の\R{量}{りょう}になったのだから、20\R{分前}{ぷん|まえ}の\R{量}{りょう}は、その\R{半分}{はん|ぶん}の $\dfrac{1}{2}$ \R{倍}{ばい}です。2-12の $2^{-1}=\dfrac{1}{2}$ と\R{合}{あ}っています。

\begin{bkcenter}
\begin{bktable}{colspec={Q[l,m,wd=16mm] *{8}{Q[c,m,wd=11mm]}}, rowsep=1.5pt, colsep=3pt, cells={font=\small}}
$x$ & $-3$ & $-2$ & $-1$ & 0 & 1 & 2 & 3 & 4 \\
$y=2^x$ & $\dfrac{1}{8}$ & $\dfrac{1}{4}$ & $\dfrac{1}{2}$ & 1 & 2 & 4 & 8 & 16 \\
\SetCell{bg=bkBlue!8} \R{小数}{しょう|すう} & 0.125 & 0.25 & 0.5 & 1 & 2 & 4 & 8 & 16 \\
\end{bktable}
\bktabcaption{$x$ が1ふえるごとに $y$ は2\R{倍}{ばい}。1へるごとに $\dfrac{1}{2}$ \R{倍}{ばい}。}
\end{bkcenter}

\noindent\begin{minipage}[c]{0.5\linewidth}\centering
\begin{tikzpicture}
\begin{axis}[bkplot, width=64mm, height=60mm, xmin=-3.6, xmax=4.6, ymin=0, ymax=17.5, xtick={-3,...,4}, ytick={4,8,12,16},
  xlabel={$x$}, ylabel={$y$}]
  \addplot[bkcurve, domain=-3.5:4.08] {2^x};
  \addplot[only marks, mark=*, mark size=1.7pt, bkInk] coordinates {(-3,0.125) (-2,0.25) (-1,0.5) (1,2) (2,4) (3,8) (4,16)};
  \addplot[only marks, mark=*, mark size=2.2pt, bkRed] coordinates {(0,1)};
  \node[font=\sffamily\footnotesize, text=bkBlue!80!black, anchor=east, fill=white, inner sep=1pt] at (axis cs:3.2,13) {$y=2^x$};
  \draw[bkRed!80!black, line width=0.5pt] (axis cs:-0.15,1.25) -- (axis cs:-1.2,5.2);
  \node[font=\sffamily\footnotesize, text=bkRed!80!black, anchor=south, fill=white, inner sep=1pt] at (axis cs:-1.4,5.4) {$(0,\ 1)$};
\end{axis}
\end{tikzpicture}
\end{minipage}\hfill
\begin{minipage}[c]{0.47\linewidth}
\R{表}{ひょう}の\R{点}{てん}をかき、なめらかな\R{線}{せん}でつなぎました（3-4）。$x$ が\R{整数}{せい|すう}でない\R{所}{ところ}の\R{値}{あたい}は、⑤で\R{決}{き}めます。
\begin{enumerate}[label={\bkcircled[bkTheme]{\arabic*}}, leftmargin=2\zw, itemsep=1pt, topsep=2pt]
\item $x=0$ で $y=1$。\R{点}{てん} $(0,\ 1)$ を\R{通}{とお}る。
\item \R{右}{みぎ}へ\R{行}{い}くほど、\R{急}{きゅう}に\R{上}{あ}がる。
\item \R{左}{ひだり}へ\R{行}{い}くと、$x$ \R{軸}{じく}に\R{近}{ちか}づくが、くっつかない。\R{何回}{なん|かい}\R{半分}{はん|ぶん}にしても、0にはならないから。
\item $y$ はいつも0より\R{大}{おお}きい。
\end{enumerate}
\end{minipage}

\par\vspace{4pt}
$y=2^x$ のように、\R{指数}{し|すう}の\R{所}{ところ}に\R{変数}{へん|すう} $x$ がある\R{関数}{かん|すう}を、\textbf{\R{指数関数}{し|すう|かん|すう}}といいます。\bkidx{しすうかんすう}{\R{指数関数}{し|すう|かん|すう}}\R{一般}{いっ|ぱん}には $y=a^x$ と\R{書}{か}きます。$a$ は1\R{回}{かい}ごとにかける\R{数}{かず}（\R{倍率}{ばい|りつ}）で、\R{正}{せい}の\R{数}{かず}です。

\textbf{$x^2$ と $2^x$ はちがう}\quad $y=x^2$ は「$x$ を2\R{回}{かい}かける」、$y=2^x$ は「2を $x$ \R{回}{かい}かける」です。くらべてみましょう。
\begin{bkcenter}
\begin{bktable}{colspec={Q[l,m,wd=16mm] *{5}{Q[c,m,wd=11mm]} Q[c,m,wd=13mm] Q[c,m,wd=20mm]}, rowsep=1.5pt, colsep=3pt, cells={font=\small}}
$x$ & 1 & 2 & 3 & 4 & 5 & 10 & 20 \\
$x^2$ & 1 & 4 & 9 & 16 & 25 & 100 & 400 \\
\SetCell{bg=bkthrlsel} $2^x$ & 2 & 4 & 8 & 16 & 32 & 1024 & 1048576 \\
\end{bktable}
\bktabcaption{$x=4$ までは\R{同}{おな}じくらい。そこから\R{先}{さき}は、$2^x$ がけたちがいに\R{大}{おお}きくなる。}
\end{bkcenter}
$x=20$ では、$x^2$ は400なのに、$2^x$ は\R{約}{やく}100\R{万}{まん}です。$x^3$ や $x^{10}$ とくらべても、$x$ を\R{大}{おお}きくしていくと、いつかはかならず $2^x$ のほうが\R{大}{おお}きくなります。

\bkthrlstep{1}{へる\R{指数関数}{し|すう|かん|すう}：$y=\left(\dfrac{1}{2}\right)^x$}
こんどは、1\R{回}{かい}ごとに\R{半分}{はん|ぶん}になる\R{量}{りょう}です。$a=\dfrac{1}{2}$ なので、$y=\left(\dfrac{1}{2}\right)^x$ です。

\begin{bkcenter}
\begin{bktable}{colspec={Q[l,m,wd=22mm] *{8}{Q[c,m,wd=11.5mm]}}, rowsep=1.5pt, colsep=3pt, cells={font=\small}}
$x$ & $-3$ & $-2$ & $-1$ & 0 & 1 & 2 & 3 & 4 \\
$y=\left(\dfrac{1}{2}\right)^x$ & 8 & 4 & 2 & 1 & 0.5 & 0.25 & 0.125 & 0.0625 \\
\end{bktable}
\end{bkcenter}

\noindent\begin{minipage}[c]{0.56\linewidth}\centering
\begin{tikzpicture}
\begin{axis}[bkplot, width=74mm, height=60mm, xmin=-4.6, xmax=4.6, ymin=0, ymax=9.6, xtick={-4,-3,-2,-1,1,2,3,4}, ytick={4,8},
  xlabel={$x$}, ylabel={$y$}]
  \addplot[bkcurve, domain=-4.3:3.2] {2^x};
  \addplot[bkcurve2, domain=-3.2:4.3] {0.5^x};
  \addplot[only marks, mark=*, mark size=1.5pt, bkInk] coordinates {(-3,8) (-2,4) (-1,2) (0,1) (1,0.5) (2,0.25) (3,0.125)};
  \node[font=\sffamily\footnotesize, text=bkBlue!80!black, anchor=west] at (axis cs:3.3,8.6) {$2^x$};
  \node[font=\sffamily\footnotesize, text=bkOrange!80!black, anchor=east] at (axis cs:-3.3,8.6) {$\left(\frac{1}{2}\right)^x$};
\end{axis}
\end{tikzpicture}
\end{minipage}\hfill
\begin{minipage}[c]{0.41\linewidth}
グラフ（オレンジ）は、\R{右}{みぎ}へ\R{行}{い}くほど\R{下}{さ}がり、$x$ \R{軸}{じく}に\R{近}{ちか}づいていきます。$2^x$ のグラフ（\R{青}{あお}）を、$y$ \R{軸}{じく}で\R{左右}{さ|ゆう}に\R{折}{お}り\R{返}{かえ}した\R{形}{かたち}です。\R{表}{ひょう}を\R{見}{み}ると、$2^x$ の\R{表}{ひょう}を\R{右}{みぎ}から\R{読}{よ}んだものになっています。
\end{minipage}

\par\vspace{4pt}
$a$ が1より\R{大}{おお}きければ、1\R{回}{かい}ごとに\R{大}{おお}きくなるので、グラフは\textbf{\R{右上}{みぎ|あ}がり}です。$a$ が0と1のあいだなら、1\R{回}{かい}ごとに\R{小}{ちい}さくなるので、グラフは\textbf{\R{右下}{みぎ|さ}がり}です。$a=1$ なら、\R{何回}{なん|かい}かけても1のままです。

\bkthrlstep{2}{「ひき\R{算}{ざん}が\R{一定}{いっ|てい}」と「わり\R{算}{ざん}が\R{一定}{いっ|てい}」}
1\R{次関数}{じ|かん|すう}と\R{指数関数}{し|すう|かん|すう}のちがいを、\R{表}{ひょう}で\R{確}{たし}かめます。$x$ を1ずつふやして、となりどうしの $y$ を、ひき\R{算}{ざん}とわり\R{算}{ざん}でくらべます。
\begin{bkcenter}
\begin{bktable}{colspec={Q[l,m,wd=38mm] *{4}{Q[c,m,wd=12mm]} Q[l,m,wd=25mm]}, rowsep=1.5pt, cells={font=\small}, column{1}={font=\sffamily\small}}
$x$ & 0 & 1 & 2 & 3 & \\
$y=2x+1$（1\R{次関数}{じ|かん|すう}） & 1 & 3 & 5 & 7 & \\
\SetCell{bg=bkBlue!10} \R{次}{つぎ}の $y$ － \R{前}{まえ}の $y$ & & 2 & 2 & 2 & \R{一定}{いっ|てい} \\
$y=2^x$（\R{指数関数}{し|すう|かん|すう}） & 1 & 2 & 4 & 8 & \\
\SetCell{bg=bkthrlsel} \R{次}{つぎ}の $y$ ÷ \R{前}{まえ}の $y$ & & 2 & 2 & 2 & \R{一定}{いっ|てい} \\
\end{bktable}
\bktabcaption{1\R{次関数}{じ|かん|すう}はひき\R{算}{ざん}が\R{一定}{いっ|てい}。\R{指数関数}{し|すう|かん|すう}はわり\R{算}{ざん}が\R{一定}{いっ|てい}。}
\end{bkcenter}
\R{指数関数}{し|すう|かん|すう}のひき\R{算}{ざん}（\R{変化}{へん|か}の\R{割合}{わり|あい}）は、1、2、4と\R{一定}{いっ|てい}ではありません。グラフがまっすぐでなく、だんだん\R{急}{きゅう}になるのはこのためです。はかったデータが\R{指数関数}{し|すう|かん|すう}かどうかは、となりどうしのわり\R{算}{ざん}が\R{一定}{いっ|てい}かで\R{見分}{み|わ}けられます。

\bksubsection{3}{\R{指数}{し|すう}のきまり}

\leavevmode\label{bk:3-11:kimari}2-12で、$x^m\times x^n=x^{m+n}$ を\R{学}{まな}びました。\R{指数関数}{し|すう|かん|すう}では、\R{指数}{し|すう}の\R{所}{ところ}が\R{変数}{へん|すう}です。そこで、\R{指数}{し|すう}のきまりを、\R{細菌}{さい|きん}の\R{時間}{じ|かん}の\R{話}{はなし}で\R{読}{よ}みなおします。ここからは、$a$ は\R{正}{せい}の\R{数}{かず}とします。

\bkthrlstep{1}{$a^{x+h}=a^x\times a^h$：\R{続}{つづ}けて\R{進}{すす}むと、\R{倍率}{ばい|りつ}はかけ\R{算}{ざん}}
\R{細菌}{さい|きん}が3\R{回}{かい}\R{分}{わ}かれて $2^3=8$ \R{倍}{ばい}になり、そのあと、さらに1\R{回}{かい}\R{分}{わ}かれたとします。\R{全部}{ぜん|ぶ}で $3+1=4$ \R{回}{かい}なので $2^{3+1}$ \R{倍}{ばい}です。いっぽう、8\R{倍}{ばい}になったものがさらに2\R{倍}{ばい}になったので、$2^3\times2^1$ \R{倍}{ばい}とも\R{言}{い}えます。
\begin{bkfigure}
\begin{tikzpicture}[x=1mm, y=1mm]
  \node[font=\small, anchor=east] at (-2,3.25) {$2^{3+1}=$};
  \foreach \i in {0,1,2} {\bkthrltape{(\i*11,0)}{9.5}{bkBlue}{2}}
  \node[font=\small] at (35,3.25) {$\times$};
  \bkthrltape{(39,0)}{9.5}{bkOrange}{2}
  \node[font=\small, anchor=west] at (51,3.25) {$=2^3\times2^1=8\times2=16$};
  \draw[bkBlue!80!black, line width=0.6pt, decorate, decoration={brace, amplitude=1.6mm, mirror}] (0,-1.8) -- (31.5,-1.8);
  \node[font=\sffamily\footnotesize, text=bkBlue!80!black, anchor=north] at (15.75,-3.9) {3\R{回}{かい}};
  \node[font=\sffamily\footnotesize, text=bkOrange!80!black, anchor=north] at (43.75,-3.9) {1\R{回}{かい}};
\end{tikzpicture}
\bkfigcaption{2を3\R{回}{かい}かけてから、さらに1\R{回}{かい}かける。\R{全部}{ぜん|ぶ}で4\R{回}{かい}。}
\end{bkfigure}
$2^4=16$ と、たしかに\R{同}{おな}じです。$x$ \R{回}{かい}\R{分}{わ}かれたあと、さらに $h$ \R{回}{かい}\R{分}{わ}かれても\R{同}{おな}じように\R{考}{かんが}えられるので、
\bkthrlline{\fbox{$a^{x+h}=a^x\times a^h$}}
これは2-12の $x^m\times x^n=x^{m+n}$ の\R{文字}{も|じ}をかえて、\R{右}{みぎ}と\R{左}{ひだり}を\R{入}{い}れかえたものです。「\R{時間}{じ|かん}を\R{続}{つづ}けて\R{進}{すす}めると、\R{倍率}{ばい|りつ}はかけ\R{算}{ざん}になる」と\R{読}{よ}めます。5-14では、このきまりを\R{使}{つか}って\R{指数関数}{し|すう|かん|すう}を\R{微分}{び|ぶん}します。

\bkthrlstep{2}{$(a^m)^n=a^{mn}$：\R{倍率}{ばい|りつ}をくり\R{返}{かえ}すと、\R{指数}{し|すう}はかけ\R{算}{ざん}}
20\R{分}{ぷん}ごとに2\R{倍}{ばい}の\R{細菌}{さい|きん}は、1\R{時間}{じ|かん}（20\R{分}{ぷん}が3\R{回}{かい}）で $2^3=8$ \R{倍}{ばい}になります。2\R{時間}{じ|かん}では、8\R{倍}{ばい}が2\R{回}{かい}なので $(2^3)^2$ \R{倍}{ばい}です。2を\R{何回}{なん|かい}かけたかを\R{数}{かぞ}えると、
\begin{bkfigure}
\begin{tikzpicture}[x=1mm, y=1mm]
  \node[font=\small, anchor=east] at (-2,3.25) {$(2^3)^2=$};
  \foreach \i in {0,1,2} {\bkthrltape{(\i*9,0)}{8}{bkBlue}{2}}
  \node[font=\small] at (28.5,3.25) {$\times$};
  \foreach \i in {0,1,2} {\bkthrltape{(32+\i*9,0)}{8}{bkBlue}{2}}
  \node[font=\small, anchor=west] at (61,3.25) {$=2^{3\times2}=2^6=64$};
  \draw[bkBlue!80!black, line width=0.6pt, decorate, decoration={brace, amplitude=1.6mm, mirror}] (0,-1.8) -- (26,-1.8);
  \draw[bkBlue!80!black, line width=0.6pt, decorate, decoration={brace, amplitude=1.6mm, mirror}] (32,-1.8) -- (58,-1.8);
  \node[font=\sffamily\footnotesize, text=bkBlue!80!black, anchor=north] at (13,-4.6) {$2^3$（1\R{時間目}{じ|かん|め}）};
  \node[font=\sffamily\footnotesize, text=bkBlue!80!black, anchor=north] at (45,-4.6) {$2^3$（2\R{時間目}{じ|かん|め}）};
\end{tikzpicture}
\bkfigcaption{3\R{回}{かい}のまとまりが2つ。2を $3\times2=6$ \R{回}{かい}かけている。}
\end{bkfigure}
$8\times8=64$ で、たしかに $2^6$ です。\R{指数}{し|すう}は「たす」のではなく「かける」ことに\R{注意}{ちゅう|い}しましょう（$2^{3+2}=32$ ではありません）。
\bkthrlline{\fbox{$(a^m)^n=a^{mn}$}}
2-7のもっと\R{知}{し}りたいの $(10^2)^3=10^6$ と\R{同}{おな}じきまりです。

\bkthrlstep{3}{$a^0=1$：0\R{回}{かい}かけると、はじめのまま}
0\R{回}{かい}\R{分}{わ}かれた（まだ\R{分}{わ}かれていない）\R{細菌}{さい|きん}は、はじめの1\R{倍}{ばい}のままです。だから $2^0=1$ です。①のきまりからも\R{確}{たし}かめられます。$a^{0+h}=a^0\times a^h$ で、\R{左}{ひだり}は $0+h=h$ なので $a^h$ です。
\[ a^h=a^0\times a^h \]
$a^h$ に $a^0$ をかけても $a^h$ のままなので、$a^0=1$ です。2-12で\R{表}{ひょう}のならびから\R{決}{き}めた $x^0=1$ と、\R{同}{おな}じ\R{答}{こた}えになりました。

\bkthrlstep{4}{$a^{-x}=\dfrac{1}{a^x}$：さかのぼると、わり\R{算}{ざん}}
$x$ \R{回}{かい}\R{進}{すす}んでから、$x$ \R{回}{かい}もどると、はじめにもどります。①のきまりで\R{計算}{けい|さん}すると、
\[ a^x\times a^{-x}=a^{x+(-x)}=a^0=1 \]
$a^{-x}$ は、$a^x$ にかけて1になる\R{数}{かず}です。だから、
\[ a^{-x}=\frac{1}{a^x} \]
です。2-12で\R{学}{まな}んだ $x^{-1}=\dfrac{1}{x}$ のつづきです。たとえば、$2^{-3}=\dfrac{1}{2^3}=\dfrac{1}{8}$ は、3\R{回前}{かい|まえ}（1\R{時間前}{じ|かん|まえ}）の\R{量}{りょう}がいまの $\dfrac{1}{8}$ だったことを\R{表}{あらわ}します。

\bksubsection{4}{\R{半分}{はん|ぶん}ずつへる：\R{半減期}{はん|げん|き}}

\leavevmode\label{bk:3-11:hangen}\R{病院}{びょう|いん}の\R{検査}{けん|さ}では、\R{弱}{よわ}い\R{放射線}{ほう|しゃ|せん}を\R{出}{だ}す\R{薬}{くすり}を\R{使}{つか}うことがあります。よく\R{使}{つか}われるテクネチウム99m という\R{物質}{ぶっ|しつ}は、6\R{時間}{じ|かん}ごとに\R{半分}{はん|ぶん}にへります。だから、\R{検査}{けん|さ}が\R{終}{お}わると\R{体}{からだ}の\R{中}{なか}からすぐに\R{少}{すく}なくなります。

\begin{bkfigure}
\begin{tikzpicture}[x=1mm, y=1mm]
  \foreach \g/\n/\hr/\pc in {0/16/0/100, 1/8/6/50, 2/4/12/25, 3/2/18/12.5, 4/1/24/6.25} {
    \begin{scope}[shift={(\g*27,0)}]
      \foreach \k in {0,...,15} {
        \pgfmathtruncatemacro\cx{mod(\k,4)}\pgfmathtruncatemacro\cy{div(\k,4)}
        \ifnum\k<\n \filldraw[fill=bkGreen!45, draw=bkGreen!70!black, line width=0.4pt] (\cx*3.4,\cy*3.4) rectangle ++(2.8,2.8);
        \else \draw[bkLine, line width=0.4pt] (\cx*3.4,\cy*3.4) rectangle ++(2.8,2.8);\fi}
      \node[font=\sffamily\footnotesize, anchor=north] at (6.4,-1.5) {\hr\R{時間}{じ|かん}};
      \node[font=\sffamily\footnotesize, anchor=south] at (6.4,14.6) {\pc\,\%};
    \end{scope}}
  \foreach \g in {0,...,3} {\draw[bk arrow, bkGray, line width=0.7pt] (\g*27+14.6,6.4) -- (\g*27+25.2,6.4);
    \node[font=\sffamily\scriptsize, anchor=south] at (\g*27+19.9,7.4) {$\times\frac{1}{2}$};}
\end{tikzpicture}
\bkfigcaption{6\R{時間}{じ|かん}ごとに、のこっている\R{量}{りょう}（\R{緑}{みどり}のマス）が\R{半分}{はん|ぶん}になる。}
\end{bkfigure}

\begin{bkcenter}
\begin{bktable}{colspec={Q[l,m,wd=38mm] *{5}{Q[c,m,wd=14mm]}}, rowsep=1.5pt, cells={font=\small}, column{1}={font=\sffamily\small}}
\R{時間}{じ|かん}（h） & 0 & 6 & 12 & 18 & 24 \\
\R{半分}{はん|ぶん}になった\R{回数}{かい|すう} & 0 & 1 & 2 & 3 & 4 \\
\SetCell{bg=bkthrlsel} のこっている\R{量}{りょう}（\%） & 100 & 50 & 25 & 12.5 & 6.25 \\
\end{bktable}
\bktabcaption{6\R{時間}{じ|かん}ごとに $\dfrac{1}{2}$ \R{倍}{ばい}。24\R{時間}{じ|かん}で $\left(\dfrac{1}{2}\right)^4=\dfrac{1}{16}$。}
\end{bkcenter}

このように、\R{量}{りょう}が\R{半分}{はん|ぶん}になるまでの\R{時間}{じ|かん}を\textbf{\R{半減期}{はん|げん|き}}といいます。\bkidx{はんげんき}{\R{半減期}{はん|げん|き}}テクネチウム99m の\R{半減期}{はん|げん|き}は6\R{時間}{じ|かん}です。100\,\%から50\,\%までも6\R{時間}{じ|かん}、50\,\%から25\,\%までも6\R{時間}{じ|かん}です。はじめの\R{量}{りょう}がいくらでも、\R{半分}{はん|ぶん}になる\R{時間}{じ|かん}は\R{同}{おな}じです。

\textbf{なぜ、いつも\R{同}{おな}じ\R{時間}{じ|かん}で\R{半分}{はん|ぶん}？}\quad この\R{物質}{ぶっ|しつ}の\R{原子}{げん|し}は、1つ1つが、\R{同}{おな}じ\R{時間}{じ|かん}のうちに\R{同}{おな}じ\R{見}{み}こみで、\R{別}{べつ}の\R{原子}{げん|し}に\R{変}{か}わります。だから、\R{原子}{げん|し}がたくさんあるほど、たくさん\R{変}{か}わります。へる\R{量}{りょう}が、いまある\R{量}{りょう}に\R{比例}{ひ|れい}するのです。\R{数学}{すう|がく}の\R{道具}{どう|ぐ}の①で\R{見}{み}たとおり、そういう\R{量}{りょう}は、\R{同}{おな}じ\R{時間}{じ|かん}ごとに\R{同}{おな}じ\R{倍率}{ばい|りつ}で\R{変}{か}わります。

\bkthrlstep{1}{\R{式}{しき}にする：$\left(\dfrac{1}{2}\right)^{\frac{t}{T}}$}
\R{半減期}{はん|げん|き}を $T$、たった\R{時間}{じ|かん}を $t$ とします。\R{半分}{はん|ぶん}になった\R{回数}{かい|すう}は、「$t$ の\R{中}{なか}に $T$ がいくつあるか」なので、$t\div T=\dfrac{t}{T}$ \R{回}{かい}です。たとえば、24\R{時間}{じ|かん}の\R{中}{なか}には6\R{時間}{じ|かん}が $\dfrac{24}{6}=4$ つあるので、4\R{回}{かい}\R{半分}{はん|ぶん}になっています。
\bkthrlline{\fbox{（のこっている\R{量}{りょう}）＝（はじめの\R{量}{りょう}）×$\left(\dfrac{1}{2}\right)^{\frac{t}{T}}$}}
$t$ と $T$ は\R{同}{おな}じ\R{単位}{たん|い}（ここでは\R{時間}{じ|かん}）にそろえます。

では、3\R{時間後}{じ|かん|ご}は\R{何}{なん}\,\%でしょう。3\R{時間}{じ|かん}は\R{半減期}{はん|げん|き}の\R{半分}{はん|ぶん}なので、
\[ \frac{t}{T}=\frac{3}{6}=0.5 \]
です。「$\dfrac{1}{2}$ を0.5\R{回}{かい}かける」とは、どういうことでしょうか。

100\,\%と50\,\%のまん\R{中}{なか}の75\,\%（\R{図}{ず}\ref{bk:3-11:fig-hangen}の\R{赤}{あか}い\R{点}{てん}）と\R{考}{かんが}えたくなります。でも、グラフは\R{直線}{ちょく|せん}ではなく、\R{曲}{ま}がっています。⑤で、\R{本当}{ほん|とう}の\R{値}{あたい}（\R{青}{あお}い\R{丸}{まる}）を\R{求}{もと}めます。

\begin{bkfigure}
\begin{tikzpicture}
\begin{axis}[bkplot, width=112mm, height=56mm, xmin=0, xmax=26.5, ymin=0, ymax=112, xtick={0,3,6,12,18,24}, ytick={25,50,75,100},
  xlabel={$t$（h）}, ylabel={\%}]
  \addplot[bkcurve, domain=0:25.5] {100*0.5^(x/6)};
  \addplot[bkRed, line width=0.9pt, dashed] coordinates {(0,100) (6,50)};
  \addplot[only marks, mark=*, mark size=1.7pt, bkInk] coordinates {(0,100) (6,50) (12,25) (18,12.5) (24,6.25)};
  \addplot[only marks, mark=*, mark size=1.9pt, bkRed] coordinates {(3,75)};
  \addplot[only marks, mark=o, mark size=2.2pt, bkBlue, line width=0.9pt] coordinates {(3,70.71)};
  \draw[bkRed!80!black, line width=0.5pt] (axis cs:3.4,76.5) -- (axis cs:8.3,93);
  \node[font=\sffamily\footnotesize, text=bkRed!80!black, anchor=west, fill=white, inner sep=1pt] at (axis cs:8.5,93) {\R{直線}{ちょく|せん}でつなぐと75\,\%？};
  \draw[bkBlue!80!black, line width=0.5pt] (axis cs:3.5,69.5) -- (axis cs:8.3,75);
  \node[font=\sffamily\footnotesize, text=bkBlue!80!black, anchor=west, fill=white, inner sep=1pt] at (axis cs:8.5,75) {\R{本当}{ほん|とう}は\R{曲線}{きょく|せん}の\R{上}{うえ}の\R{値}{あたい}};
\end{axis}
\end{tikzpicture}
\bkfigcaption{\R{半減期}{はん|げん|き}6\R{時間}{じ|かん}の\R{物質}{ぶっ|しつ}。3\R{時間後}{じ|かん|ご}は、\R{赤}{あか}い\R{点}{てん}（75\,\%）ではなく、\R{曲線}{きょく|せん}の\R{上}{うえ}の\R{青}{あお}い\R{丸}{まる}。}\label{bk:3-11:fig-hangen}
\end{bkfigure}

\bksubsection{5}{\R{分数}{ぶん|すう}の\R{指数}{し|すう}：$x^{\frac{1}{2}}=\sqrt{x}$}

\leavevmode\label{bk:3-11:bunsu}\R{指数}{し|すう}が\R{整数}{せい|すう}なら、「\R{何回}{なん|かい}かけるか」でした。0.5\R{回}{かい}かけることはできません。そこで、これまでのきまりがそのまま\R{成}{な}り\R{立}{た}つように、\R{分数}{ぶん|すう}の\R{指数}{し|すう}の\R{意味}{い|み}を\R{決}{き}めます（2-7で、0\R{乗}{じょう}やマイナスの\R{指数}{し|すう}を\R{決}{き}めたのと\R{同}{おな}じやり\R{方}{かた}です）。ここでは2\R{通}{とお}りで\R{考}{かんが}えます。

\bkthrlstep{1}{\R{表}{ひょう}で\R{考}{かんが}える}
\R{細菌}{さい|きん}の\R{数}{かず} $y=2^x$ を、$x$ を0.5ずつ\R{進}{すす}めた\R{表}{ひょう}にします。$x=0.5$ は「20\R{分}{ぷん}の\R{半分}{はん|ぶん}の10\R{分}{ぷん}たったとき」です。
\begin{bkfigure}
\begin{tikzpicture}[x=1mm, y=1mm]
  \foreach \i in {0,...,4} \draw[bkLine, line width=0.6pt] (30+\i*18,-6) -- (30+\i*18,6);
  \draw[bkLine, line width=0.6pt] (0,6) rectangle (120,-6) (0,0) -- (120,0);
  \fill[bkTheme!12] (0.3,5.7) rectangle (29.7,0.3); \fill[bkTheme!12] (0.3,-0.3) rectangle (29.7,-5.7);
  \node[font=\sffamily\small, anchor=west] at (2,3) {$x$};
  \node[font=\sffamily\small, anchor=west] at (2,-3) {$y=2^x$};
  \foreach \v [count=\i from 0] in {0,0.5,1,1.5,2} \node[font=\small] at (39+\i*18,3) {$\v$};
  \foreach \v [count=\i from 0] in {1,{\color{bkRed}？},2,{\color{bkRed}？},4} \node[font=\small] at (39+\i*18,-3) {$\v$};
  \foreach \i in {0,...,3} {
    \pgfmathsetmacro\xa{39+\i*18}\pgfmathsetmacro\xb{39+(\i+1)*18}
    \draw[-{Stealth[length=4pt]}, bkOrange, line width=0.9pt, shorten <=1mm, shorten >=1mm] (\xa,-6.8) .. controls (\xa+4,-11.5) and (\xb-4,-11.5) .. (\xb,-6.8);
    \node[font=\sffamily\footnotesize, text=bkOrange!80!black, anchor=north, inner sep=0.5pt] at ({(\xa+\xb)/2},-10.8) {$\times r$};}
  \foreach \i in {0,1} {
    \pgfmathsetmacro\xa{39+\i*36}\pgfmathsetmacro\xb{39+(\i+1)*36}
    \draw[-{Stealth[length=4pt]}, bkBlue, line width=0.9pt, shorten <=1mm, shorten >=1mm] (\xa,6.8) .. controls (\xa+6,13) and (\xb-6,13) .. (\xb,6.8);
    \node[font=\sffamily\footnotesize, text=bkBlue!80!black, anchor=south, inner sep=0.5pt] at ({(\xa+\xb)/2},11.8) {$\times2$};}
\end{tikzpicture}
\bkfigcaption{$x$ が0.5\R{進}{すす}むごとに、\R{同}{おな}じ\R{倍率}{ばい|りつ} $r$ をかける。$r$ を2\R{回}{かい}かけると2\R{倍}{ばい}。}
\end{bkfigure}
\R{指数関数}{し|すう|かん|すう}は、$x$ が\R{同}{おな}じだけ\R{進}{すす}むと、いつも\R{同}{おな}じ\R{倍率}{ばい|りつ}になるのでした（②）。$x$ が0.5\R{進}{すす}むときの\R{倍率}{ばい|りつ}を $r$ とします。
\begin{enumerate}[label={\bkcircled[bkTheme]{\arabic*}}, leftmargin=2.4\zw, itemsep=1pt, topsep=2pt]
\item $x=0$ から $x=1$ へは、0.5ずつ2\R{回}{かい}\R{進}{すす}むので、$r$ を2\R{回}{かい}かけます。
\item いっぽう、$x=0$ から1へは2\R{倍}{ばい}です。だから $r\times r=2$。
\item 2\R{乗}{じょう}すると2になる\R{正}{せい}の\R{数}{かず}は $\sqrt{2}=1.414\cdots$（2-8）。だから $r=\sqrt{2}$ です。
\end{enumerate}
$x=0.5$ の\R{値}{あたい}は $1\times r=\sqrt{2}$ なので、$2^{0.5}=2^{\frac{1}{2}}=\sqrt{2}$ です。10\R{分後}{ぷん|ご}には、\R{細菌}{さい|きん}は\R{約}{やく}1.41\R{倍}{ばい}になっています。1.5\R{倍}{ばい}ではありません。

\bkthrlstep{2}{\R{指数}{し|すう}のきまりで\R{考}{かんが}える}
③の②のきまり $(a^m)^n=a^{mn}$ が、\R{分数}{ぶん|すう}の\R{指数}{し|すう}でも\R{成}{な}り\R{立}{た}つとします。$x^{\frac{1}{2}}$ を2\R{乗}{じょう}すると、
\[ \Bigl(x^{\frac{1}{2}}\Bigr)^2=x^{\frac{1}{2}\times2}=x^1=x \]
$x^{\frac{1}{2}}$ は「2\R{乗}{じょう}すると $x$ になる\R{数}{かず}」です。それは $\sqrt{x}$ です（2-8。\R{正}{せい}の\R{数}{かず}のほうにします。$a^x$ はいつも\R{正}{せい}だからです）。③の①の $a^{x+h}=a^x\times a^h$ でも、$x^{\frac{1}{2}}\times x^{\frac{1}{2}}=x^{\frac{1}{2}+\frac{1}{2}}=x^1=x$ で、\R{同}{おな}じことがわかります。
\bkthrlline{\fbox{$x^{\frac{1}{2}}=\sqrt{x}$}\qquad \fbox{$x^{-\frac{1}{2}}=\dfrac{1}{x^{\frac{1}{2}}}=\dfrac{1}{\sqrt{x}}$}}
2つめの\R{式}{しき}は、③の④の $a^{-x}=\dfrac{1}{a^x}$ を\R{使}{つか}いました。たとえば、
\[ 9^{\frac{1}{2}}=\sqrt{9}=3,\qquad 9^{-\frac{1}{2}}=\frac{1}{\sqrt{9}}=\frac{1}{3},\qquad 2^{\frac{1}{2}}=\sqrt{2}=1.414\cdots \]
\R{表}{ひょう}で\R{考}{かんが}えても、きまりで\R{考}{かんが}えても、\R{同}{おな}じ\R{答}{こた}えになりました。\R{電卓}{でん|たく}の \texttt{x\^{}y} のキー（\R{表計算}{ひょう|けい|さん}なら \texttt{=2\^{}0.5}）で $2^{0.5}$ を\R{計算}{けい|さん}すると、$\sqrt{2}$ のキーと\R{同}{おな}じ1.414213562が\R{出}{で}ます。

\noindent\begin{minipage}[t]{0.42\linewidth}\vspace{0pt}
$\sqrt{x}$ を $x^{\frac{1}{2}}$ と\R{書}{か}けると、$\sqrt{\phantom{x}}$ も「\R{指数}{し|すう}のなかま」になります。5-6・5-8では、この\R{書}{か}き\R{方}{かた}で $\sqrt{x}$ を\R{微分}{び|ぶん}します。
\end{minipage}\hfill
\begin{minipage}[t]{67mm}\vspace{0pt}\bkkigocard{$x^{\frac{1}{2}}$}{エックスの2\R{分}{ぶん}の1\R{乗}{じょう}}{$\sqrt{x}$ のこと。\R{分数}{ぶん|すう}の\R{指数}{し|すう}}{x to the one-half}{3-11}\end{minipage}
\par\vspace{6pt}\bkidx{ぶんすうのしすう}{\R{分数}{ぶん|すう}の\R{指数}{し|すう}}

\bkthrlstep{3}{\R{半減期}{はん|げん|き}の\R{半分}{はん|ぶん}の\R{時間}{じ|かん}では}
④の3\R{時間後}{じ|かん|ご}にもどります。のこっている\R{量}{りょう}は
\[ 100\times\left(\frac{1}{2}\right)^{0.5}=100\times\frac{1}{2^{0.5}}=100\times\frac{1}{\sqrt{2}}=\frac{100}{1.414}=70.7 \]
で、\R{約}{やく}70.7\,\%です。75\,\%ではありません。\R{確}{たし}かめてみましょう。3\R{時間}{じ|かん}ごとに0.707\R{倍}{ばい}なら、6\R{時間}{じ|かん}では $0.707\times0.707=0.4998\cdots$ で、ちゃんと\R{約}{やく}\R{半分}{はん|ぶん}になります。もし75\,\%だとすると、6\R{時間}{じ|かん}で $0.75\times0.75=0.5625$ \R{倍}{ばい}になってしまい、\R{半分}{はん|ぶん}より\R{多}{おお}くのこります。\R{図}{ず}\ref{bk:3-11:fig-hangen}で、\R{青}{あお}い\R{丸}{まる}が\R{曲線}{きょく|せん}の\R{上}{うえ}にあることを\R{確}{たし}かめましょう。

\begin{motto}[3]
\par\addvspace{3pt}\textbf{3\R{乗}{じょう}するともとにもどる\R{数}{かず}：$x^{\frac{1}{3}}$}\par\noindent
\R{同}{おな}じように\R{考}{かんが}えると、$\Bigl(x^{\frac{1}{3}}\Bigr)^3=x^{\frac{1}{3}\times3}=x$ なので、$x^{\frac{1}{3}}$ は「3\R{乗}{じょう}すると $x$ になる\R{数}{かず}」です。たとえば $2\times2\times2=8$ なので $8^{\frac{1}{3}}=2$、$10\times10\times10=1000$ なので $1000^{\frac{1}{3}}=10$ です。\R{体積}{たい|せき}が\qty{1000}{cm^3}（\qty{1}{L}）の\R{立方体}{りっ|ぽう|たい}の1\R{辺}{ぺん}は、$1000^{\frac{1}{3}}=10$ で\qty{10}{cm}です。
\end{motto}

\bksubsection{6}{\R{整数}{せい|すう}でない\R{指数}{し|すう}を、\R{電卓}{でん|たく}で\R{計算}{けい|さん}する}

\leavevmode\label{bk:3-11:dentaku}$2^{1.5}$ を\R{求}{もと}めます。1.5は $1+0.5$ なので、③の①のきまりで2つに\R{分}{わ}けられます。
\[ 2^{1.5}=2^{1+0.5}=2^1\times2^{0.5}=2\times1.414=2.828 \]
\R{電卓}{でん|たく}で「2 \texttt{x\^{}y} 1.5 ＝」と\R{押}{お}すと、2.828427… が\R{出}{で}ます。\R{同}{おな}じ\R{答}{こた}えです。

\noindent\begin{minipage}[c]{0.66\linewidth}
$2^{1.4}$ はどうでしょう。1.4は $\dfrac{14}{10}$ なので、③の②のきまりで
\[ \bigl(2^{1.4}\bigr)^{10}=2^{1.4\times10}=2^{14}=16384 \]
です。$2^{1.4}$ は「10\R{乗}{じょう}すると16384になる\R{数}{かず}」です。\R{手}{て}で\R{求}{もと}めるのはたいへんですが、\R{電卓}{でん|たく}なら「2 \texttt{x\^{}y} 1.4 ＝」で2.639015… とすぐに\R{出}{で}ます。\R{表計算}{ひょう|けい|さん}では \texttt{=2\^{}1.4} と\R{書}{か}きます。
\end{minipage}\hfill
\begin{minipage}[c]{0.3\linewidth}\centering
\begin{tikzpicture}[x=1mm, y=1mm]
  \fill[bkGray!70, rounded corners=2mm] (0,0) rectangle (38,30);
  \bkdisplay{(2.5,17)}{33}{2.639016}
  \foreach \i in {0,...,3} {\fill[white, rounded corners=0.8mm] (3.5+\i*8.2,3.5) rectangle ++(6.3,5);}
  \fill[bkOrange!60, rounded corners=0.8mm] (3.5,10) rectangle ++(14.5,5);
  \node[font=\ttfamily\scriptsize] at (10.75,12.5) {x\^{}y};
\end{tikzpicture}
\end{minipage}

\par\vspace{4pt}
こうして、$x$ がどんな\R{数}{かず}でも $2^x$ の\R{値}{あたい}が\R{決}{き}まります。②のグラフの\R{点}{てん}と\R{点}{てん}のあいだを、なめらかな\R{線}{せん}でつないでよいのは、このためです。

\begin{seikaku}
\R{指数}{し|すう}が\R{分数}{ぶん|すう}でない\R{数}{かず}（$\sqrt{2}$ や $\pi$ のように、\R{小数}{しょう|すう}がどこまでも\R{続}{つづ}く\R{数}{かず}。2-9）のときは、その\R{数}{かず}にかぎりなく\R{近}{ちか}い\R{分数}{ぶん|すう}の\R{指数}{し|すう}で\R{計算}{けい|さん}した\R{値}{あたい}が\R{近}{ちか}づく\R{先}{さき}として、$a^x$ を\R{決}{き}めます。また、$a$ がマイナスだと、$(-2)^{\frac{1}{2}}$（2\R{乗}{じょう}して $-2$ になる\R{数}{かず}）のように\R{決}{き}められない\R{値}{あたい}が\R{出}{で}てくるので、\R{指数関数}{し|すう|かん|すう} $a^x$ では $a$ を\R{正}{せい}の\R{数}{かず}にします。
\end{seikaku}

\begin{column}{\R{米}{こめ}1\R{粒}{つぶ}から\R{始}{はじ}める「\R{倍々}{ばい|ばい}のほうび」}
\noindent\begin{minipage}[c]{0.7\linewidth}
\R{豊臣秀吉}{とよ|とみ|ひで|よし}に\R{仕}{つか}えた\R{曽呂利新左衛門}{そ|ろ|り|しん|ざ|え|もん}は、ほうびに「1\R{日目}{にち|め}は\R{米}{こめ}1\R{粒}{つぶ}、\R{2日}{ふつか}\R{目}{め}は2\R{粒}{つぶ}、\R{3日}{みっか}\R{目}{め}は4\R{粒}{つぶ}と、\R{前}{まえ}の\R{日}{ひ}の2\R{倍}{ばい}ずつ、30\R{日間}{にち|かん}ください」とたのんだ、という\R{話}{はなし}が\R{伝}{つた}わっています。1\R{日目}{にち|め}が $2^0=1$ \R{粒}{つぶ}なので、$n$ \R{日目}{にち|め}は $2^{n-1}$ \R{粒}{つぶ}です。30\R{日目}{にち|め}の1\R{日分}{にち|ぶん}だけで $2^{29}=536870912$ \R{粒}{つぶ}。\R{米}{こめ}は1\,kg で\R{約}{やく}5\R{万}{まん}\R{粒}{つぶ}なので、$536870912\div50000$ で\R{約}{やく}10700\,kg、つまり\R{約}{やく}11トンにもなります。
\end{minipage}\hfill
\begin{minipage}[c]{0.27\linewidth}\centering
\begin{tikzpicture}[x=1.2mm, y=1.2mm]
  \foreach \k/\n in {0/1,1/2,2/4,3/8} {
    \foreach \j in {1,...,\n} {\pgfmathsetmacro\cx{\k*7.5+mod(\j-1,2)*2.6}\pgfmathsetmacro\cy{div(\j-1,2)*2.2}
      \fill[bkYellow!25] (\cx,\cy) ellipse (1.1 and 0.75); \draw[bkYellow!60!black, line width=0.3pt] (\cx,\cy) ellipse (1.1 and 0.75);}
    \node[font=\sffamily\scriptsize, anchor=north] at (\k*7.5+1.3,-1.5) {\n};}
  \node[font=\sffamily\scriptsize, anchor=north] at (13.5,-5.5) {1〜4\R{日目}{にち|め}（\R{粒}{つぶ}）};
\end{tikzpicture}
\end{minipage}
\end{column}

\begin{taisetsu}
ふえる（へる）\R{量}{りょう}が、いまある\R{量}{りょう}に\R{比例}{ひ|れい}するとき：\par
\hspace*{1\zw}（$x$ \R{回後}{かい|ご}の\R{量}{りょう}）＝（はじめの\R{量}{りょう}）×$a^x$（$a$ は1\R{回}{かい}ごとの\R{倍率}{ばい|りつ}）\par\vspace{2pt}
$y=a^x$：$a>1$ ならふえる、$0<a<1$ ならへる。\R{点}{てん} $(0,\ 1)$ を\R{通}{とお}り、$y$ はいつも\R{正}{せい}。\par\vspace{2pt}
$a^{x+h}=a^x\times a^h$、\ $(a^m)^n=a^{mn}$、\ $a^0=1$、\ $a^{-x}=\dfrac{1}{a^x}$\par\vspace{2pt}
$x^{\frac{1}{2}}=\sqrt{x}$、\ $x^{-\frac{1}{2}}=\dfrac{1}{\sqrt{x}}$\par\vspace{2pt}
\R{半減期}{はん|げん|き} $T$：（のこっている\R{量}{りょう}）＝（はじめの\R{量}{りょう}）×$\left(\dfrac{1}{2}\right)^{\frac{t}{T}}$
\end{taisetsu}

\begin{kotoba}
\small
\bkword{\R{指数関数}{し|すう|かん|すう}}{$y=a^x$ のように、\R{指数}{し|すう}の\R{所}{ところ}に\R{変数}{へん|すう}がある\R{関数}{かん|すう}。\R{倍々}{ばい|ばい}でふえる・へる\R{量}{りょう}を\R{表}{あらわ}す。}
\bkword{\R{半減期}{はん|げん|き}}{\R{量}{りょう}が\R{半分}{はん|ぶん}になるまでの\R{時間}{じ|かん}。はじめの\R{量}{りょう}がいくらでも\R{同}{おな}じ。}
\bkword{\R{分数}{ぶん|すう}の\R{指数}{し|すう}}{$x^{\frac{1}{2}}=\sqrt{x}$ のように、\R{指数}{し|すう}のきまりが\R{成}{な}り\R{立}{た}つように\R{決}{き}めた\R{累乗}{るい|じょう}。}
\end{kotoba}

%% ---------------- 例題 ----------------
\clearpage
\bksection{reidai}{\R{例題}{れい|だい}でたしかめよう}\bkidx{しすうのきまり}{\R{指数}{し|すう}のきまり}

\begin{reidai}[\R{倍々}{ばい|ばい}でふえる\R{量}{りょう}]
30\R{分}{ぷん}ごとに2\R{倍}{ばい}にふえる\R{細菌}{さい|きん}が、はじめに3000\R{個}{こ}います。\par
（1）3\R{時間後}{じ|かん|ご}には\R{何個}{なん|こ}になりますか。\par
（2）はじめて100\R{万}{まん}\R{個}{こ}をこえるのは、\R{何時間何分後}{なん|じ|かん|なん|ぷん|ご}ですか。\par
\kangaekata 30\R{分}{ぷん}が\R{何回}{なん|かい}あるかを\R{数}{かぞ}えて、2をその\R{回数}{かい|すう}だけかけます（①）。（2）は\R{表}{ひょう}を\R{作}{つく}って、100\R{万}{まん}をこえる\R{所}{ところ}をさがします。\par
\shiki（1）\nobreak $\qty{3}{h}\times\qty{60}{min/h}=\qty{180}{min}$。$180\div30=6$ \R{回}{かい}\par
\hspace*{2.6\zw}$3000\times2^6=3000\times64=192000$\par
\hspace*{2.6\zw}（2）\nobreak 6\R{回}{かい}で192000。\R{続}{つづ}けて2\R{倍}{ばい}していく\par
\hspace*{2.6\zw}7\R{回}{かい}：384000、\ 8\R{回}{かい}：768000、\ 9\R{回}{かい}：1536000\par
\hspace*{2.6\zw}9\R{回目}{かい|め}ではじめてこえる。$30\times9=270$ で\qty{270}{min}\par
\hspace*{2.6\zw}$\qty{270}{min}\div\qty{60}{min/h}=\qty{4.5}{h}$。0.5\R{時間}{じ|かん}は30\R{分}{ぷん}\par
\kotae（1）\nobreak 192000\R{個}{こ}\quad（2）\nobreak 4\R{時間}{じ|かん}30\R{分後}{ぷん|ご}
\end{reidai}

\begin{jibun}
20\R{分}{ぷん}ごとに2\R{倍}{ばい}にふえる\R{細菌}{さい|きん}が、はじめに500\R{個}{こ}います。（1）2\R{時間後}{じ|かん|ご}には\R{何個}{なん|こ}になりますか。（2）はじめて10\R{万}{まん}\R{個}{こ}をこえるのは、\R{何時間何分後}{なん|じ|かん|なん|ぷん|ご}ですか。
\bkjans{（1）\nobreak \qty{120}{min}÷\qty{20}{min}＝6\R{回}{かい}。$500\times2^6=32000$ で32000\R{個}{こ}。（2）\nobreak 7\R{回}{かい}で64000、8\R{回}{かい}で128000。$20\times8=160$ で\qty{160}{min}、2\R{時間}{じ|かん}40\R{分後}{ぷん|ご}}
\end{jibun}

\begin{bkcenter}
\begin{bktable}{colspec={Q[l,m,wd=12mm] *{12}{Q[c,m,wd=7.2mm]}}, rowsep=1.2pt, colsep=1.5pt, cells={font=\footnotesize}, row{1}={bg=bkTheme!10}}
$n$ & 1 & 2 & 3 & 4 & 5 & 6 & 7 & 8 & 9 & 10 & 11 & 12 \\
$2^n$ & 2 & 4 & 8 & 16 & 32 & 64 & 128 & 256 & 512 & 1024 & 2048 & 4096 \\
\end{bktable}
\bktabcaption{よく\R{使}{つか}う2の\R{累乗}{るい|じょう}。\R{右}{みぎ}へ1つ\R{進}{すす}むごとに2\R{倍}{ばい}。$2^{10}=1024$ は、ほぼ1000。}
\end{bkcenter}

\begin{reidai}[\R{割合}{わり|あい}でふえる・へる]
（1）1\R{年}{ねん}ごとに、いまある\R{金額}{きん|がく}の10\,\%の\R{利息}{り|そく}がつく\R{貯金}{ちょ|きん}に、10000\R{円}{えん}を\R{預}{あず}けました。3\R{年後}{ねん|ご}はいくらですか。\par
（2）200\R{万円}{まん|えん}の\R{車}{くるま}の\R{値段}{ね|だん}が、1\R{年}{ねん}ごとに、そのときの\R{値段}{ね|だん}の20\,\%ずつ\R{下}{さ}がります。2\R{年後}{ねん|ご}は\R{何万円}{なん|まん|えん}ですか。\par
\kangaekata 10\,\%ふえる＝1＋0.1＝1.1\R{倍}{ばい}、20\,\%へる＝1－0.2＝0.8\R{倍}{ばい}です（2-5）。それを\R{年数}{ねん|すう}だけかけます（①）。\par
\shiki（1）\nobreak $1.1^3=1.1\times1.1\times1.1=1.21\times1.1=1.331$\par
\hspace*{2.6\zw}$10000\times1.331=13310$\par
\hspace*{2.6\zw}（2）\nobreak $0.8^2=0.8\times0.8=0.64$\par
\hspace*{2.6\zw}$200\times0.64=128$\par
\kotae（1）\nobreak 13310\R{円}{えん}\quad（2）\nobreak 128\R{万円}{まん|えん}
\end{reidai}
\noindent（1）で「\R{毎年}{まい|とし}1000\R{円}{えん}ずつ」と\R{考}{かんが}えると13000\R{円}{えん}です。ちがいの310\R{円}{えん}は、\R{利息}{り|そく}についた\R{利息}{り|そく}です。

\begin{jibun}
（1）1\R{年}{ねん}ごとに、いまある\R{金額}{きん|がく}の5\,\%の\R{利息}{り|そく}がつく\R{貯金}{ちょ|きん}に、40000\R{円}{えん}を\R{預}{あず}けました。2\R{年後}{ねん|ご}はいくらですか。（2）10\R{万円}{まん|えん}のパソコンの\R{値段}{ね|だん}が、1\R{年}{ねん}ごとに、そのときの\R{値段}{ね|だん}の10\,\%ずつ\R{下}{さ}がります。3\R{年後}{ねん|ご}は\R{何円}{なん|えん}ですか。
\bkjans{（1）\nobreak $1.05^2=1.1025$、$40000\times1.1025=44100$ で44100\R{円}{えん}。（2）\nobreak $0.9^3=0.729$、$100000\times0.729=72900$ で72900\R{円}{えん}}
\end{jibun}

\begin{reidai}[\R{指数}{し|すう}のきまりを\R{使}{つか}う]
\R{次}{つぎ}の\R{数}{かず}を\R{求}{もと}めましょう。（5）は $2^x$ を\R{使}{つか}って\R{表}{あらわ}しましょう。\par
（1）\nobreak $3^{2+1}$ を $3^2\times3^1$ に\R{分}{わ}けて\quad（2）\nobreak $(2^2)^3$\quad（3）\nobreak $5^0$\quad（4）\nobreak $2^{-3}$\quad（5）\nobreak $2^{x+3}$\par
\kangaekata ③の4つのきまりを\R{使}{つか}います。（2）は\R{指数}{し|すう}をかけ\R{算}{ざん}、（5）は $a^{x+h}=a^x\times a^h$ で $h=3$ です。\par
\shiki（1）\nobreak $3^{2+1}=3^2\times3^1=9\times3=27$\quad（\R{確}{たし}かめ：$3^3=27$）\par
\hspace*{2.6\zw}（2）\nobreak $(2^2)^3=2^{2\times3}=2^6=64$\quad（\R{確}{たし}かめ：$4^3=64$）\par
\hspace*{2.6\zw}（3）\nobreak $5^0=1$\quad（4）\nobreak $2^{-3}=\dfrac{1}{2^3}=\dfrac{1}{8}$\par
\hspace*{2.6\zw}（5）\nobreak $2^{x+3}=2^x\times2^3=8\times2^x$\par
\kotae（1）\nobreak 27\quad（2）\nobreak 64\quad（3）\nobreak 1\quad（4）\nobreak $\dfrac{1}{8}$\quad（5）\nobreak $8\times2^x$
\end{reidai}
\noindent（5）は、「$x$ \R{回目}{かい|め}から、さらに3\R{回}{かい}\R{進}{すす}むと、いつでも8\R{倍}{ばい}になる」という\R{意味}{い|み}です。20\R{分}{ぷん}ごとに2\R{倍}{ばい}の\R{細菌}{さい|きん}なら、いつから\R{数}{かぞ}えても、1\R{時間後}{じ|かん|ご}には8\R{倍}{ばい}です。

\begin{jibun}
\R{次}{つぎ}の\R{数}{かず}を\R{求}{もと}めましょう。（5）は $3^x$ を\R{使}{つか}って\R{表}{あらわ}しましょう。（1）\nobreak $2^{4+1}$ を $2^4\times2^1$ に\R{分}{わ}けて（2）\nobreak $(3^2)^2$（3）\nobreak $10^0$（4）\nobreak $10^{-2}$（5）\nobreak $3^{x+2}$
\bkjans{（1）\nobreak $16\times2=32$（2）\nobreak $3^4=81$（3）\nobreak 1（4）\nobreak $\dfrac{1}{100}$（5）\nobreak $3^x\times3^2=9\times3^x$}
\end{jibun}

\begin{reidai}[\R{半減期}{はん|げん|き}]
セシウム137という\R{物質}{ぶっ|しつ}の\R{半減期}{はん|げん|き}は30\R{年}{ねん}です。\par
（1）90\R{年後}{ねん|ご}には、はじめの\R{何分}{なん|ぶん}の1になりますか。\par
（2）はじめの $\dfrac{1}{16}$ になるのは、\R{何年後}{なん|ねん|ご}ですか。\par
\kangaekata \R{半分}{はん|ぶん}になる\R{回数}{かい|すう}は $\dfrac{t}{T}$（④）。（2）は、$\dfrac{1}{2}$ を\R{何回}{なん|かい}かけると $\dfrac{1}{16}$ になるかを\R{考}{かんが}えます。\par
\shiki（1）\nobreak $\dfrac{90}{30}=3$ \R{回}{かい}。$\left(\dfrac{1}{2}\right)^3=\dfrac{1}{2\times2\times2}=\dfrac{1}{8}$\par
\hspace*{2.6\zw}（2）\nobreak $2\times2\times2\times2=16$ なので、$\dfrac{1}{16}=\left(\dfrac{1}{2}\right)^4$。4\R{回}{かい}\par
\hspace*{2.6\zw}$30\times4=120$\par
\kotae（1）\nobreak $\dfrac{1}{8}$\quad（2）\nobreak 120\R{年後}{ねん|ご}
\end{reidai}

\begin{jibun}
テクネチウム99m の\R{半減期}{はん|げん|き}は6\R{時間}{じ|かん}です。（1）30\R{時間後}{じ|かん|ご}には、はじめの\R{何分}{なん|ぶん}の1になりますか。（2）はじめの $\dfrac{1}{8}$ になるのは、\R{何時間後}{なん|じ|かん|ご}ですか。
\bkjans{（1）\nobreak $\dfrac{30}{6}=5$ \R{回}{かい}。$\left(\dfrac{1}{2}\right)^5=\dfrac{1}{32}$（2）\nobreak $\dfrac{1}{8}=\left(\dfrac{1}{2}\right)^3$ で3\R{回}{かい}。$6\times3=18$ で18\R{時間後}{じ|かん|ご}}
\end{jibun}

\begin{bkfigure}
\begin{tikzpicture}[x=1mm, y=1mm]
  \foreach \k/\f in {0/1,1/{\frac{1}{2}},2/{\frac{1}{4}},3/{\frac{1}{8}},4/{\frac{1}{16}}} {
    \pgfmathsetmacro\w{64/2^\k}
    \fill[bkGreen!40] (\k*26,0) rectangle ++(\w/3,8);
    \draw[bkGreen!70!black, line width=0.5pt] (\k*26,0) rectangle ++(64/3,8);
    \node[font=\small, anchor=south] at (\k*26+10.67,9.5) {$\displaystyle\f$};}
  \foreach \k/\y in {0/0,1/30,2/60,3/90,4/120} {\node[font=\sffamily\footnotesize, anchor=north] at (\k*26+10.67,-1.5) {\y\R{年後}{ねん|ご}};}
\end{tikzpicture}
\bkfigcaption{セシウム137は、30\R{年}{ねん}ごとに\R{半分}{はん|ぶん}。\R{何回}{なん|かい}\R{半分}{はん|ぶん}になったかは $\dfrac{t}{T}$ で\R{数}{かぞ}える。}
\end{bkfigure}

\begin{reidai}[\R{分数}{ぶん|すう}の\R{指数}{し|すう}]
\R{次}{つぎ}の\R{数}{かず}を\R{求}{もと}めましょう。\par
（1）\nobreak $9^{\frac{1}{2}}$\qquad（2）\nobreak $16^{-\frac{1}{2}}$\qquad（3）\nobreak $4^{1.5}$\qquad（4）\nobreak $100^{\frac{1}{2}}$\par
\kangaekata $x^{\frac{1}{2}}=\sqrt{x}$、$x^{-\frac{1}{2}}=\dfrac{1}{\sqrt{x}}$（⑤）。（3）は $1.5=1+0.5$ と\R{分}{わ}けて、$a^{x+h}=a^x\times a^h$ を\R{使}{つか}います。\par
\shiki（1）\nobreak $9^{\frac{1}{2}}=\sqrt{9}=3$\quad 3を2\R{乗}{じょう}すると9\par
\hspace*{2.6\zw}（2）\nobreak $16^{-\frac{1}{2}}=\dfrac{1}{\sqrt{16}}=\dfrac{1}{4}$\par
\hspace*{2.6\zw}（3）\nobreak $4^{1.5}=4^{1+0.5}=4^1\times4^{0.5}=4\times\sqrt{4}=4\times2=8$\par
\hspace*{2.6\zw}（4）\nobreak $100^{\frac{1}{2}}=\sqrt{100}=10$\par
\kotae（1）\nobreak 3\quad（2）\nobreak $\dfrac{1}{4}$\quad（3）\nobreak 8\quad（4）\nobreak 10
\end{reidai}

\begin{jibun}
\R{次}{つぎ}の\R{数}{かず}を\R{求}{もと}めましょう。（1）\nobreak $25^{\frac{1}{2}}$\quad（2）\nobreak $9^{-\frac{1}{2}}$\quad（3）\nobreak $9^{1.5}$\quad（4）\nobreak $10000^{\frac{1}{2}}$
\bkjans{（1）\nobreak 5（2）\nobreak $\dfrac{1}{3}$（3）\nobreak $9\times\sqrt{9}=9\times3=27$（4）\nobreak 100}
\end{jibun}

\noindent\bkreidairef{5}の（3）を、⑤の\R{表}{ひょう}の\R{考}{かんが}え\R{方}{かた}でも\R{確}{たし}かめましょう。$4^x$ は、$x$ が1\R{進}{すす}むと4\R{倍}{ばい}です。$x$ が0.5\R{進}{すす}むときの\R{倍率}{ばい|りつ}は、2\R{回}{かい}かけて4になる\R{数}{かず}なので、$\sqrt{4}=2$ です。
\begin{bkfigure}
\begin{tikzpicture}[x=1mm, y=1mm]
  \foreach \i in {0,...,4} \draw[bkLine, line width=0.6pt] (24+\i*18,-6) -- (24+\i*18,6);
  \draw[bkLine, line width=0.6pt] (0,6) rectangle (114,-6) (0,0) -- (114,0);
  \fill[bkTheme!12] (0.3,5.7) rectangle (23.7,0.3); \fill[bkTheme!12] (0.3,-0.3) rectangle (23.7,-5.7);
  \node[font=\small, anchor=west] at (2,3) {$x$};
  \node[font=\small, anchor=west] at (2,-3) {$4^x$};
  \foreach \v [count=\i from 0] in {0,0.5,1,1.5,2} \node[font=\small] at (33+\i*18,3) {$\v$};
  \foreach \v [count=\i from 0] in {1,2,4,8,16} \node[font=\small] at (33+\i*18,-3) {$\v$};
  \foreach \i in {0,...,3} {
    \pgfmathsetmacro\xa{33+\i*18}\pgfmathsetmacro\xb{33+(\i+1)*18}
    \draw[-{Stealth[length=4pt]}, bkOrange, line width=0.9pt, shorten <=1mm, shorten >=1mm] (\xa,-6.8) .. controls (\xa+4,-11.5) and (\xb-4,-11.5) .. (\xb,-6.8);
    \node[font=\sffamily\footnotesize, text=bkOrange!80!black, anchor=north, inner sep=0.5pt] at ({(\xa+\xb)/2},-10.8) {$\times2$};}
\end{tikzpicture}
\bkfigcaption{$x$ が0.5\R{進}{すす}むごとに2\R{倍}{ばい}。$x=1.5$ では $4^{1.5}=8$ で、\R{例題}{れい|だい}5の（3）と\R{同}{おな}じ。}
\end{bkfigure}

\begin{reidai}[\R{半減期}{はん|げん|き}の\R{途中}{と|ちゅう}]
\R{半減期}{はん|げん|き}が6\R{時間}{じ|かん}の\R{物質}{ぶっ|しつ}が、はじめに400\,mg あります。（1）9\R{時間後}{じ|かん|ご}と、（2）15\R{時間後}{じ|かん|ご}には、\R{約}{やく}\R{何}{なん}mg のこっていますか。$\left(\dfrac{1}{2}\right)^{0.5}=\dfrac{1}{\sqrt{2}}=0.707$ を\R{使}{つか}いましょう。\par
\kangaekata \R{半分}{はん|ぶん}になる\R{回数}{かい|すう} $\dfrac{t}{T}$ が1.5や2.5になります。\R{整数}{せい|すう}の\R{部分}{ぶ|ぶん}と0.5に\R{分}{わ}けます（③の①）。\par
\shiki（1）\nobreak $\dfrac{9}{6}=1.5$。$\left(\dfrac{1}{2}\right)^{1.5}=\left(\dfrac{1}{2}\right)^1\times\left(\dfrac{1}{2}\right)^{0.5}=0.5\times0.707$\par
\hspace*{2.6\zw}$400\times0.5\times0.707=200\times0.707=141.4$\par
\hspace*{2.6\zw}（2）\nobreak $\dfrac{15}{6}=2.5$。$\left(\dfrac{1}{2}\right)^{2.5}=\left(\dfrac{1}{2}\right)^2\times\left(\dfrac{1}{2}\right)^{0.5}=0.25\times0.707$\par
\hspace*{2.6\zw}$400\times0.25\times0.707=100\times0.707=70.7$\par
\kotae（1）\nobreak \R{約}{やく}141\,mg\quad（2）\nobreak \R{約}{やく}70.7\,mg
\end{reidai}
\noindent（1）の\R{答}{こた}えは、6\R{時間後}{じ|かん|ご}の200\,mg と12\R{時間後}{じ|かん|ご}の100\,mg のあいだにあります。まん\R{中}{なか}の150\,mg より\R{少}{すく}し\R{少}{すく}ないのは、グラフが\R{下}{した}にふくらんだ\R{曲線}{きょく|せん}だからです。

\begin{jibun}
セシウム137（\R{半減期}{はん|げん|き}30\R{年}{ねん}）が、はじめに2000\,g あります。（1）15\R{年後}{ねん|ご}と、（2）45\R{年後}{ねん|ご}には、\R{約}{やく}\R{何}{なん}g のこっていますか。$\left(\dfrac{1}{2}\right)^{0.5}=0.707$ を\R{使}{つか}いましょう。
\bkjans{（1）\nobreak $\dfrac{15}{30}=0.5$。$2000\times0.707=1414$ で\R{約}{やく}1414\,g（2）\nobreak $\dfrac{45}{30}=1.5$。$2000\times0.5\times0.707=707$ で\R{約}{やく}707\,g}
\end{jibun}

\begin{reidai}[\R{整数}{せい|すう}でない\R{指数}{し|すう}と\R{電卓}{でん|たく}]
（1）$3^{1.5}$ を、③のきまりと $\sqrt{3}=1.732$ を\R{使}{つか}って\R{求}{もと}め、\R{電卓}{でん|たく}で\R{確}{たし}かめましょう。\par
（2）\R{空気}{くう|き}を、\R{熱}{ねつ}がにげないように\R{急}{きゅう}におしちぢめると、\R{圧力}{あつ|りょく} $p$ と\R{体積}{たい|せき} $V$ のあいだに「$pV^{1.4}$ が\R{一定}{いっ|てい}」という\R{関係}{かん|けい}があります。はじめの\R{圧力}{あつ|りょく}と\R{体積}{たい|せき}をどちらも1として（はじめの\R{何倍}{なん|ばい}かで\R{表}{あらわ}す）、\R{体積}{たい|せき}を\R{急}{きゅう}に0.5にすると、\R{圧力}{あつ|りょく}は\R{何倍}{なん|ばい}になりますか。\R{電卓}{でん|たく}で $0.5^{1.4}=0.379$ を\R{求}{もと}めて\R{使}{つか}いましょう。\par
\kangaekata （2）はじめの $pV^{1.4}$ は、$1\times1^{1.4}=1$ です（1は\R{何回}{なん|かい}かけても1）。おしちぢめたあとも、$pV^{1.4}$ は1のままです。\par
\shiki（1）\nobreak $3^{1.5}=3^1\times3^{0.5}=3\times\sqrt{3}=3\times1.732=5.196$\par
\hspace*{2.6\zw}\R{電卓}{でん|たく}：「3 \texttt{x\^{}y} 1.5 ＝」で5.196152…\par
\hspace*{2.6\zw}（2）\nobreak $p\times0.5^{1.4}=1$ より $p\times0.379=1$\par
\hspace*{2.6\zw}\R{両辺}{りょう|へん}を0.379でわって（2-14）、$p=\dfrac{1}{0.379}=2.638\cdots$\par
\kotae（1）\nobreak \R{約}{やく}5.196\quad（2）\nobreak \R{約}{やく}2.64\R{倍}{ばい}
\end{reidai}
\noindent ゆっくりおしちぢめて\R{温度}{おん|ど}が\R{変}{か}わらないときは、$pV$ が\R{一定}{いっ|てい}（3-7のボイルの\R{法則}{ほう|そく}）なので、\R{体積}{たい|せき}を\R{半分}{はん|ぶん}にすると\R{圧力}{あつ|りょく}は2\R{倍}{ばい}です。\R{急}{きゅう}におしちぢめると、\R{空気}{くう|き}の\R{温度}{おん|ど}が\R{上}{あ}がるので、\R{圧力}{あつ|りょく}は2\R{倍}{ばい}より\R{大}{おお}きくなります。

\begin{jibun}
（1）$5^{1.5}$ を、③のきまりと $\sqrt{5}=2.236$ を\R{使}{つか}って\R{求}{もと}めましょう。（2）\R{例題}{れい|だい}の（2）と\R{同}{おな}じ\R{空気}{くう|き}を、\R{急}{きゅう}に\R{体積}{たい|せき}0.25までおしちぢめると、\R{圧力}{あつ|りょく}は\R{約}{やく}\R{何倍}{なん|ばい}になりますか。\R{電卓}{でん|たく}で $0.25^{1.4}=0.1436$ を\R{求}{もと}めて\R{使}{つか}いましょう。
\bkjans{（1）\nobreak $5\times\sqrt{5}=5\times2.236=11.18$（2）\nobreak $p\times0.1436=1$ より $p=\dfrac{1}{0.1436}=6.96\cdots$ で\R{約}{やく}6.96\R{倍}{ばい}}
\end{jibun}

%% ---------------- まちがえやすいところ ----------------
\bksection{tsumazuki}{まちがえやすいところをチェック}\label{bk:3-11:tsuma}

\begin{tsumazuki}{かけ\R{算}{ざん}なのに、\R{指数}{し|すう}をかける}
\bkng{$2^3\times2^2=2^{3\times2}=2^6$}
\bkok{$2^3\times2^2=2^{3+2}=2^5=32$}
2を3\R{回}{かい}かけて、さらに2\R{回}{かい}かけるので、\R{全部}{ぜん|ぶ}で5\R{回}{かい}です（$8\times4=32$）。\R{指数}{し|すう}をかけるのは $(2^3)^2$ のときです。
\end{tsumazuki}

\begin{tsumazuki}{かっこの\R{累乗}{るい|じょう}なのに、\R{指数}{し|すう}をたす}
\bkng{$(2^3)^2=2^{3+2}=2^5$}
\bkok{$(2^3)^2=2^{3\times2}=2^6=64$}
$2^3=8$ を2\R{回}{かい}かけるので $8\times8=64$。3\R{回}{かい}のまとまりが2つで、2を6\R{回}{かい}かけています。
\end{tsumazuki}

\begin{tsumazuki}{10\,\%ずつ3\R{年}{ねん}で、30\,\%ふえる？}
\bkng{\R{毎年}{まい|とし}10\,\%ずつふえるので、3\R{年}{ねん}で30\,\%ふえて1.3\R{倍}{ばい}}
\bkok{$1.1^3=1.331$ で、\R{約}{やく}33\,\%ふえる}
2\R{年目}{ねん|め}からは、ふえた\R{分}{ぶん}にも10\,\%がつきます。たし\R{算}{ざん}ではなく、1.1のかけ\R{算}{ざん}をくり\R{返}{かえ}します。
\end{tsumazuki}

\begin{tsumazuki}{\R{半減期}{はん|げん|き}の\R{半分}{はん|ぶん}の\R{時間}{じ|かん}で、$\dfrac{3}{4}$ のこる？}
\bkng{3\R{時間後}{じ|かん|ご}（\R{半減期}{はん|げん|き}6\R{時間}{じ|かん}の\R{半分}{はん|ぶん}）は、100\,\%と50\,\%のまん\R{中}{なか}で75\,\%}
\bkok{$\left(\dfrac{1}{2}\right)^{0.5}=\dfrac{1}{\sqrt{2}}=0.707$ で、\R{約}{やく}70.7\,\%}
75\,\%が2\R{回}{かい}だと $0.75\times0.75=0.5625$ で、\R{半分}{はん|ぶん}になりません。グラフは\R{直線}{ちょく|せん}ではなく\R{曲線}{きょく|せん}です。
\end{tsumazuki}

\begin{tsumazuki}{2\R{分}{ぶん}の1\R{乗}{じょう}を、2でわる}
\bkng{$9^{\frac{1}{2}}=9\div2=4.5$}
\bkok{$9^{\frac{1}{2}}=\sqrt{9}=3$}
$x^{\frac{1}{2}}$ は「2\R{乗}{じょう}すると $x$ になる\R{数}{かず}」です。$4.5\times4.5=20.25$ で、9になりません。$4^{\frac{1}{2}}=2$ は $4\div2$ と\R{偶然}{ぐう|ぜん}\R{同}{おな}じなので、\R{気}{き}をつけましょう。
\end{tsumazuki}

\begin{tsumazuki}{マイナスの\R{指数}{し|すう}で、マイナスの\R{数}{かず}になる？}
\bkng{$2^{-3}=-8$}
\bkok{$2^{-3}=\dfrac{1}{2^3}=\dfrac{1}{8}$}
\R{指数}{し|すう}のマイナスは「わる」ことを\R{表}{あらわ}します（2-7）。\R{指数関数}{し|すう|かん|すう}の\R{値}{あたい}は、いつも\R{正}{せい}の\R{数}{かず}です（②のグラフ）。
\end{tsumazuki}

%% ---------------- 物理・工学 ----------------
\begin{butsuri}
\par\addvspace{3pt}\textbf{\R{冷}{さ}めるお\R{湯}{ゆ}：\R{部屋}{へ|や}との\R{温度}{おん|ど}の\R{差}{さ}が、\R{同}{おな}じ\R{割合}{わり|あい}でへる}\par\vspace{4pt}% 第14回 A3：となりの囲みとそろえる
\noindent\begin{minipage}[c]{0.66\linewidth}
80\,℃のお\R{湯}{ゆ}を、20\,℃の\R{部屋}{へ|や}に\R{置}{お}きます。お\R{湯}{ゆ}は、\R{部屋}{へ|や}との\R{温度}{おん|ど}の\R{差}{さ}が\R{大}{おお}きいほど\R{速}{はや}く\R{冷}{さ}めます。\R{冷}{さ}める\R{速}{はや}さ（へる\R{量}{りょう}）が、いまある\R{温度}{おん|ど}の\R{差}{さ}に\R{比例}{ひ|れい}するのです。だから、\R{温度}{おん|ど}の\R{差}{さ}は、\R{同}{おな}じ\R{時間}{じ|かん}ごとに\R{同}{おな}じ\R{倍率}{ばい|りつ}でへります。あるマグカップでは、\R{差}{さ}が1\R{分}{ぷん}ごとに0.95\R{倍}{ばい}になりました。
\end{minipage}\hfill
\begin{minipage}[c]{0.3\linewidth}\centering
\begin{tikzpicture}[x=1mm, y=1mm]
  \filldraw[fill=bkRed!12, draw=bkInk!80, line width=0.9pt, rounded corners=1.2mm] (0,0) rectangle (18,20);
  \fill[bkOrange!35] (1,1) rectangle (17,16.5);
  \draw[bkInk!80, line width=1.6pt] (18,15) arc (90:-90:4.5);
  \foreach \x in {5,9,13} {\draw[bkGray, line width=0.7pt] (\x,22) .. controls (\x-2,25) and (\x+2,27) .. (\x,30);}
  \node[font=\sffamily\footnotesize, anchor=north] at (9,-1.5) {80\,℃→？};
\end{tikzpicture}
\end{minipage}
\par\vspace{2pt}
\begin{bkcenter}
\begin{bktable}{colspec={Q[l,m,wd=34mm] *{7}{Q[c,m,wd=10.4mm]}}, rowsep=1pt, colsep=2pt, cells={font=\footnotesize}, column{1}={font=\sffamily\footnotesize}}
\R{時間}{じ|かん}（\R{分}{ふん}） & 0 & 1 & 2 & 3 & 5 & 10 & 20 \\
\R{温度}{おん|ど}の\R{差}{さ}（℃） & 60 & 57 & 54.2 & 51.4 & 46.4 & 35.9 & 21.5 \\
\SetCell{bg=bkthrlsel} お\R{湯}{ゆ}の\R{温度}{おん|ど}（℃） & 80 & 77 & 74.2 & 71.4 & 66.4 & 55.9 & 41.5 \\
\end{bktable}
\bktabcaption{\R{差}{さ}は $60\times0.95^{t}$。お\R{湯}{ゆ}の\R{温度}{おん|ど}は、それに\R{部屋}{へ|や}の20\,℃をたしたもの。}
\end{bkcenter}
\noindent\begin{minipage}[c]{0.5\linewidth}\centering
\begin{tikzpicture}
\begin{axis}[bkplot, width=62mm, height=54mm, xmin=0, xmax=33, ymin=0, ymax=92, xtick={0,10,20,30}, ytick={20,40,60,80},
  xlabel={\R{分}{ふん}}, ylabel={℃}]
  \addplot[bkthin, dashed, domain=0:31] {20};
  \addplot[bkcurve2, domain=0:31] {20+60*0.95^x};
  \addplot[only marks, mark=*, mark size=1.5pt, bkInk] coordinates {(0,80) (10,55.92) (20,41.51) (30,32.88)};
  \node[font=\sffamily\scriptsize, anchor=north east, fill=bkPurpleBg, inner sep=1pt] at (axis cs:31,17.5) {\R{部屋}{へ|や}の20\,℃};
\end{axis}
\end{tikzpicture}
\end{minipage}\hfill
\begin{minipage}[c]{0.47\linewidth}
はじめの1\R{分}{ぷん}で3\,℃\R{下}{さ}がりますが、10\R{分後}{ぷん|ご}からの1\R{分}{ぷん}では、\R{約}{やく}1.8\,℃しか\R{下}{さ}がりません。\R{差}{さ}が\R{小}{ちい}さくなるほど、ゆっくり\R{冷}{さ}めるのです。\R{温度}{おん|ど}を1\R{分}{ぷん}ごとのとびとびでなく、いつでも\R{使}{つか}える1つの\R{式}{しき}で\R{書}{か}くには、5-14で\R{学}{まな}ぶふしぎな\R{数}{かず}を\R{使}{つか}います（5-18・6-4）。
\end{minipage}
\end{butsuri}

\begin{butsuri}
\par\addvspace{3pt}\textbf{ディーゼルエンジン：\R{火花}{ひ|ばな}なしで\R{燃}{も}える}\par\vspace{7pt}% 第14回 A3：見出しは1行で本文の幅に収め、下に少し余白（1行目にルビがないので、となりの囲みと見た目の間をそろえて 7pt）
\noindent\begin{minipage}[c]{0.68\linewidth}
トラックなどのディーゼルエンジンは、シリンダー（つつ）の\R{中}{なか}の\R{空気}{くう|き}を、ピストンで\R{体積}{たい|せき}が\R{約}{やく} $\dfrac{1}{16}$ になるまで\R{急}{きゅう}におしちぢめます。\bkreidairef{7}と\R{同}{おな}じように「$pV^{1.4}$ が\R{一定}{いっ|てい}」で\R{計算}{けい|さん}します。\R{電卓}{でん|たく}で $\left(\dfrac{1}{16}\right)^{1.4}=0.0625^{1.4}=0.0206$ なので、
\bkthrlline{$p\times0.0206=1$\quad より\quad $p=\dfrac{1}{0.0206}=48.5\cdots$}
\R{圧力}{あつ|りょく}は\R{約}{やく}49\R{倍}{ばい}になります（\R{実際}{じっ|さい}は\R{熱}{ねつ}が\R{少}{すこ}しにげるので、これより\R{小}{ちい}さめ）。ゆっくりなら16\R{倍}{ばい}です。\R{温度}{おん|ど}も\R{同}{おな}じような\R{計算}{けい|さん}（くわしくは\R{物理}{ぶつ|り}シリーズ『\R{熱学}{ねつ|がく}』）で\R{求}{もと}められ、\R{約}{やく}640\,℃になります。\R{実際}{じっ|さい}は\R{熱}{ねつ}がにげるので\R{少}{すこ}し\R{低}{ひく}めですが、\R{軽油}{けい|ゆ}がひとりでに\R{燃}{も}えだす\R{温度}{おん|ど}（\R{約}{やく}250\,℃）よりずっと\R{高}{たか}い\R{温度}{おん|ど}です。そこに\R{燃料}{ねん|りょう}の\R{軽油}{けい|ゆ}をふきこむと、\R{火花}{ひ|ばな}がなくても\R{燃}{も}えだします。
\end{minipage}\hfill
\begin{minipage}[c]{0.28\linewidth}\centering
\begin{tikzpicture}[x=1mm, y=1mm]
  \fill[bkGray!35] (0,0) rectangle (2.5,36) (17.5,0) rectangle (20,36);
  \fill[bkGray!35] (0,36) rectangle (20,38.5);
  \fill[bkRed!18] (2.5,26) rectangle (17.5,36);
  \filldraw[fill=bkGray!60, draw=bkInk!70, line width=0.6pt] (2.8,20) rectangle (17.2,26);
  \fill[bkGray!75] (8.5,4) rectangle (11.5,20);
  \draw[bk arrow, bkBlue, line width=1pt] (14.5,8) -- (14.5,17);
  \fill[bkInk!80] (8.5,38.5) rectangle (11.5,42);
  \node[font=\sffamily\scriptsize] at (10,31) {\R{空気}{くう|き}};
  \draw[bkInk!70, line width=0.4pt] (15,23) -- (21.3,23);
  \node[font=\sffamily\scriptsize, anchor=west] at (21.5,23) {ピストン};
  \node[font=\sffamily\scriptsize, anchor=west] at (13,41) {\R{燃料}{ねん|りょう}};
\end{tikzpicture}
\end{minipage}

\par\addvspace{6pt}\textbf{\R{大学}{だい|がく}への\R{種}{たね}}\par\vspace{3pt}\noindent
\R{細菌}{さい|きん}・お\R{湯}{ゆ}・\R{放射線}{ほう|しゃ|せん}・\R{電気}{でん|き}の\R{回路}{かい|ろ}（コンデンサーにためた\R{電気}{でん|き}を\R{抵抗}{てい|こう}に\R{流}{なが}すと、\R{電圧}{でん|あつ}は\R{同}{おな}じ\R{時間}{じ|かん}ごとに\R{同}{おな}じ\R{割合}{わり|あい}でへる）は、どれも「\R{変}{か}わる\R{速}{はや}さが、いまある\R{量}{りょう}に\R{比例}{ひ|れい}する」しくみです。この「しくみ」を\R{式}{しき}にしたものが\R{微分方程式}{び|ぶん|ほう|てい|しき}で、その\R{答}{こた}えが\R{指数関数}{し|すう|かん|すう}になります（5-14・6-4）。
\end{butsuri}

\begin{column}{X\R{線}{せん}\R{写真}{しゃ|しん}で、\R{骨}{ほね}が\R{白}{しろ}くうつるわけ}
\noindent\begin{minipage}[c]{0.66\linewidth}
X\R{線}{せん}は、\R{物}{もの}の\R{中}{なか}を\R{進}{すす}むと、\R{同}{おな}じ\R{厚}{あつ}さを\R{通}{とお}るごとに\R{同}{おな}じ\R{割合}{わり|あい}で\R{弱}{よわ}くなります。\R{水}{みず}に\R{近}{ちか}い\R{筋肉}{きん|にく}などでは、\R{約}{やく}3〜4\,cm \R{進}{すす}むごとにおよそ\R{半分}{はん|ぶん}になります。\R{骨}{ほね}では、\R{約}{やく}1\,cm でおよそ\R{半分}{はん|ぶん}です。たとえば\R{骨}{ほね}を\qty{4}{cm}\R{通}{とお}ると、$\left(\dfrac{1}{2}\right)^4=\dfrac{1}{16}$ まで\R{弱}{よわ}くなります。\R{骨}{ほね}の\R{後}{うし}ろにはX\R{線}{せん}があまり\R{届}{とど}かないので、フィルムが\R{黒}{くろ}くならず、\R{白}{しろ}くうつるのです。
\end{minipage}\hfill
\begin{minipage}[c]{0.3\linewidth}\centering
\begin{tikzpicture}[x=1mm, y=1mm]
  \fill[bkInk!85, rounded corners=1.5mm] (0,0) rectangle (26,36);
  \fill[bkGray!55] (6,0) .. controls (4,12) and (4.5,26) .. (6.5,36) -- (19.5,36) .. controls (21.5,26) and (22,12) .. (20,0) -- cycle;
  \fill[white, rounded corners=1.5mm] (11,4) rectangle (15,32);
  \fill[white] (13,4) ellipse (4 and 2.3) (13,32) ellipse (4 and 2.3);
  \fill[bkGray!55] (13,1.9) circle (1.1) (13,34.1) circle (1.1);
  \draw[bkInk!60, line width=0.4pt] (19,18) -- (29,18); \node[font=\sffamily\scriptsize, anchor=west] at (29.3,18) {\R{筋肉}{きん|にく}};
  \draw[bkInk!60, line width=0.4pt] (14.6,26) -- (29,26); \node[font=\sffamily\scriptsize, anchor=west] at (29.3,26) {\R{骨}{ほね}};
  \node[font=\sffamily\scriptsize, anchor=north] at (13,-1.5) {X\R{線}{せん}\R{写真}{しゃ|しん}（\R{足}{あし}）};
\end{tikzpicture}
\end{minipage}
\end{column}

%% ---------------- 章末問題 ----------------
\bkendproblems
\bkproblemgroup{\R{基本}{き|ほん}}{1}

\bkproblem 20\R{分}{ぷん}ごとに2\R{倍}{ばい}にふえる\R{細菌}{さい|きん}が、はじめに100\R{個}{こ}います。（1）2\R{時間後}{じ|かん|ご}には\R{何個}{なん|こ}になりますか。（2）はじめて1\R{万}{まん}\R{個}{こ}をこえるのは、\R{何時間何分後}{なん|じ|かん|なん|ぷん|ご}ですか。
\bkpans{（1）\nobreak 6\R{回}{かい}。$100\times2^6=6400$ で6400\R{個}{こ}。（2）\nobreak 6\R{回}{かい}で6400、7\R{回}{かい}で12800。$20\times7=140$ で\qty{140}{min}、2\R{時間}{じ|かん}20\R{分後}{ぷん|ご}}

\bkproblem （1）1\R{年}{ねん}ごとに、いまある\R{金額}{きん|がく}の4\,\%の\R{利息}{り|そく}がつく\R{貯金}{ちょ|きん}に、50000\R{円}{えん}を\R{預}{あず}けました。2\R{年後}{ねん|ご}はいくらですか。（2）80\R{万円}{まん|えん}の\R{機械}{き|かい}の\R{値打}{ね|う}ちが、1\R{年}{ねん}ごとに、そのときの25\,\%ずつ\R{下}{さ}がります。2\R{年後}{ねん|ご}は\R{何万円}{なん|まん|えん}ですか。
\bkpans{（1）\nobreak $1.04^2=1.0816$、$50000\times1.0816=54080$ で54080\R{円}{えん}。（2）\nobreak $0.75^2=0.5625$、$80\times0.5625=45$ で45\R{万円}{まん|えん}}

\bkproblem \R{次}{つぎ}の\R{数}{かず}を\R{求}{もと}めましょう。（5）は $5^x$ を\R{使}{つか}って\R{表}{あらわ}しましょう。\par
\hspace*{1\zw}（1）\nobreak $2^{5+1}$ を $2^5\times2^1$ に\R{分}{わ}けて\quad（2）\nobreak $(2^3)^3$\quad（3）\nobreak $8^0$\quad（4）\nobreak $3^{-2}$\quad（5）\nobreak $5^{x+1}$
\bkpans{（1）\nobreak $32\times2=64$（2）\nobreak $2^9=512$（3）\nobreak 1（4）\nobreak $\dfrac{1}{9}$（5）\nobreak $5\times5^x$}

\bkproblem \R{炭素}{たん|そ}14の\R{半減期}{はん|げん|き}は\R{約}{やく}5730\R{年}{ねん}です。（1）17190\R{年後}{ねん|ご}には、はじめの\R{何分}{なん|ぶん}の1になりますか。（2）はじめの $\dfrac{1}{4}$ になるのは\R{何年後}{なん|ねん|ご}ですか。
\bkpans{（1）\nobreak $\dfrac{17190}{5730}=3$ \R{回}{かい}。$\dfrac{1}{8}$（2）\nobreak 2\R{回}{かい}。$5730\times2=11460$ で11460\R{年後}{ねん|ご}}

\bkproblem \R{次}{つぎ}の\R{数}{かず}を\R{求}{もと}めましょう。（1）\nobreak $36^{\frac{1}{2}}$\quad（2）\nobreak $25^{-\frac{1}{2}}$\quad（3）\nobreak $4^{2.5}$\quad（4）\nobreak $16^{1.5}$
\bkpans{（1）\nobreak 6（2）\nobreak $\dfrac{1}{5}$（3）\nobreak $4^2\times4^{0.5}=16\times2=32$（4）\nobreak $16\times\sqrt{16}=16\times4=64$}

\bkproblem \R{半減期}{はん|げん|き}が6\R{時間}{じ|かん}の\R{物質}{ぶっ|しつ}が、はじめに800\,mg あります。（1）3\R{時間後}{じ|かん|ご}と、（2）9\R{時間後}{じ|かん|ご}には、\R{約}{やく}\R{何}{なん}mg のこっていますか。$\left(\dfrac{1}{2}\right)^{0.5}=0.707$ を\R{使}{つか}いましょう。
\bkpans{（1）\nobreak 0.5\R{回}{かい}。$800\times0.707=565.6$ で\R{約}{やく}566\,mg（2）\nobreak 1.5\R{回}{かい}。$800\times0.5\times0.707=282.8$ で\R{約}{やく}283\,mg}

\bkproblemgroup{チャレンジ}{2}

\bkproblem 1\R{時間}{じ|かん}ごとに2\R{倍}{ばい}にふえる\R{細菌}{さい|きん}が、はじめに1000\R{個}{こ}います。1\R{時間}{じ|かん}30\R{分後}{ぷん|ご}には\R{約}{やく}\R{何個}{なん|こ}になりますか。$\sqrt{2}=1.414$ を\R{使}{つか}い、\R{一}{いち}の\R{位}{くらい}まで\R{答}{こた}えましょう。
\bkpans{1\R{時間}{じ|かん}30\R{分}{ぷん}は1.5\R{回}{かい}。$2^{1.5}=2\times1.414=2.828$、$1000\times2.828=2828$ で\R{約}{やく}2828\R{個}{こ}}

\bkproblem ガソリンエンジンでは、シリンダーの\R{空気}{くう|き}（とガソリンの\R{霧}{きり}）を、\R{体積}{たい|せき}が\R{約}{やく} $\dfrac{1}{10}$ になるまで\R{急}{きゅう}におしちぢめます。$pV^{1.4}$ が\R{一定}{いっ|てい}として、はじめの\R{圧力}{あつ|りょく}と\R{体積}{たい|せき}を1とすると、\R{圧力}{あつ|りょく}は\R{約}{やく}\R{何倍}{なん|ばい}になりますか。\R{電卓}{でん|たく}で $0.1^{1.4}$ を\R{求}{もと}めて、\R{小数第}{しょう|すう|だい}1\R{位}{い}まで\R{答}{こた}えましょう。
\bkpans{$0.1^{1.4}=0.0398$。$p\times0.0398=1$ より $p=\dfrac{1}{0.0398}=25.1\cdots$ で\R{約}{やく}25.1\R{倍}{ばい}}

\bkproblem ある\R{町}{まち}の\R{人口}{じん|こう}は、1\R{年}{ねん}ごとに、そのときの\R{人口}{じん|こう}の2\,\%ずつへっています。10\R{年後}{ねん|ご}の\R{人口}{じん|こう}は、いまの\R{約}{やく}\R{何}{なん}\,\%ですか。\R{電卓}{でん|たく}を\R{使}{つか}い、\R{小数第}{しょう|すう|だい}1\R{位}{い}まで\R{答}{こた}えましょう（20\,\%へると\R{考}{かんが}えるのは、まちがいです）。
\bkpans{1\R{年}{ねん}ごとに0.98\R{倍}{ばい}。$0.98^{10}=0.8170\cdots$ で\R{約}{やく}81.7\,\%}

\begin{furikaeri}
\begin{checklist}
\item 「ふえる\R{量}{りょう}が、いまある\R{量}{りょう}に\R{比例}{ひ|れい}する」\R{変化}{へん|か}を、\R{表}{ひょう}・\R{式}{しき}・グラフにできる。\bkfurref{bk:3-11:hirei}
\item \R{指数}{し|すう}のきまり $a^{x+h}=a^x\times a^h$、$(a^m)^n=a^{mn}$、$a^0=1$ を\R{使}{つか}える。\bkfurref{bk:3-11:kimari}
\item $x^{\frac{1}{2}}=\sqrt{x}$ のわけを\R{説明}{せつ|めい}できる。\bkfurref{bk:3-11:bunsu}
\end{checklist}
\end{furikaeri}

\bkfukidashi{delta}{\R{倍々}{ばい|ばい}は、はじめはゆっくりでも、あとで\R{一気}{いっ|き}に\R{大}{おお}きくなるね。\R{次}{つぎ}の\R{章}{しょう}では\R{反対}{はん|たい}に、「\R{何回}{なん|かい}かけたら、その\R{数}{かず}になるか」を\R{求}{もと}める\R{対数}{たい|すう}を\R{学}{まな}ぶよ。}
